%% =====================================================================
%%  Quantum Kernels for EEG: An Invariance Mismatch Masquerading as Advantage
%%
%%  Target journal : Journal of Neural Engineering (IOP, Q1)
%%  Document class : iopjournal.cls (IOP Publishing, official package)
%%                    Fetch with: bash paper/get_iop_class.sh
%%
%%  Build:   latexmk -pdf main.tex
%%  Tables and inline numbers come from tables_auto.tex, which is
%%  regenerated from the result CSVs by:  python paper/make_tables.py
%%  Never edit numbers by hand.
%% =====================================================================
\documentclass{iopjournal}

\usepackage{graphicx}
%% iopjournal has no quarrel with amsmath, unlike the legacy iopart
%% class, so these can be loaded directly.
\usepackage{amsmath,amssymb}
\usepackage[colorlinks=true,allcolors=blue]{hyperref}

%% Prefer the paper-mode figures: no in-figure title or note (the \caption
%% supplies both) and a white canvas. Generate with:
%%     PYTHONPATH=src python -m qeeg.figures --paper
%%     PYTHONPATH=src python -m qeeg.figures_reference --paper
\graphicspath{{../results/figures_paper/}{../results/figures/}{figures/}}

%% Auto-generated \newcommand macros for every number quoted in the prose.
%% Table floats live in tables_auto.tex and are input in the document body,
%% since a float is not legal in the preamble.
%% AUTO-GENERATED by paper/make_tables.py: DO NOT EDIT BY HAND.
%% Regenerate after every benchmark run:  python paper/make_tables.py


%% -------------------------------------------- inline macros

\newcommand{\NSubjects}{30}
\newcommand{\NPipelines}{15}
\newcommand{\NFoldScores}{6\,750}
\newcommand{\NTrials}{45}
\newcommand{\BestClassicalName}{classical/CSP+LDA}
\newcommand{\BestClassicalAcc}{0.611}
\newcommand{\BestQuantumName}{quantum/CNOT-kernel-SVM}
\newcommand{\BestQuantumAcc}{0.579}
\newcommand{\PrimaryDelta}{+0.032}
\newcommand{\PrimaryP}{\ensuremath{{}={}}0.002}
\newcommand{\PrimaryDz}{0.58}
\newcommand{\PrimaryBetter}{24/30}
\newcommand{\AblationDelta}{+0.0106}
\newcommand{\AblationP}{\ensuremath{{}={}}0.309}
\newcommand{\AblationBetter}{17/30}
\newcommand{\DimDelta}{+0.0163}
\newcommand{\DimP}{\ensuremath{{}={}}0.173}
\newcommand{\NSurviveHolm}{0}
\newcommand{\NFamilyTests}{14}
\newcommand{\FastestName}{classical/logvar+LDA}
\newcommand{\FastestSec}{0.046}
\newcommand{\SlowestName}{quantum/Bures-RBF-SVM}
\newcommand{\SlowestSec}{12.4}
\newcommand{\SpeedRatio}{272}

\newcommand{\ConcNSubjects}{10}

%% -------------------------------- cross-dataset macros

\newcommand{\BciNSubjects}{9}
\newcommand{\BciNTrials}{288}
\newcommand{\BciNFoldScores}{2\,025}
\newcommand{\BciBestClassicalAcc}{0.763}
\newcommand{\BciBestQuantumAcc}{0.685}
\newcommand{\BciWorstQuantumAcc}{0.580}
\newcommand{\GapPhysio}{0.032}
\newcommand{\GapBci}{0.079}
\newcommand{\GapWorstPhysio}{0.088}
\newcommand{\GapWorstBci}{0.183}
\newcommand{\SpearmanRho}{0.882}
\newcommand{\BciPrimaryDelta}{+0.066}
\newcommand{\BciPrimaryP}{\ensuremath{{}={}}0.004}
\newcommand{\BciPrimaryDz}{1.14}
\newcommand{\BciPrimaryBetter}{9/9}
\newcommand{\BciDensDelta}{+0.171}
\newcommand{\BciDensDz}{1.45}
\newcommand{\BciDensBetter}{9/9}
\newcommand{\FisherPrimaryP}{\ensuremath{{}<{}}0.001}
\newcommand{\FisherDensP}{\ensuremath{{}<{}}0.001}
\newcommand{\FisherAblationP}{\ensuremath{{}={}}0.237}
\newcommand{\BciWilcoxonFloor}{0.0039}
\newcommand{\BciHolmFloor}{0.0547}


%% ------------------------------ reference-frame macros

\newcommand{\NPipelinesExt}{23}
\newcommand{\NPipelinesFb}{15}
\newcommand{\PhysFrameMin}{+0.050}
\newcommand{\PhysFrameMax}{+0.106}
\newcommand{\PhysFrameBestKernel}{Fidelity}
\newcommand{\PhysFrameMaxP}{0.049}
\newcommand{\PhysFrameAllBetter}{no}
\newcommand{\PhysFrameN}{30}
\newcommand{\TwinName}{control/riemann-kernel-SVM}
\newcommand{\TwinAcc}{0.621}
\newcommand{\TwinDeltaMin}{-0.0151}
\newcommand{\TwinDeltaMax}{+0.0081}
\newcommand{\TwinMinP}{\ensuremath{{}={}}0.183}
\newcommand{\HeadBestClassical}{classical/CSP+LDA}
\newcommand{\HeadSensorDelta}{+0.0523}
\newcommand{\HeadSensorP}{\ensuremath{{}={}}0.016}
\newcommand{\HeadSensorBetter}{20/30}
\newcommand{\HeadRefDelta}{-0.0183}
\newcommand{\HeadRefLead}{0.0183}
\newcommand{\HeadRefP}{\ensuremath{{}={}}0.059}
\newcommand{\HeadRefBetter}{10/30}
\newcommand{\HeadRefKernel}{quantum/Fidelity-ref-SVM}
\newcommand{\BciFrameMin}{+0.164}
\newcommand{\BciFrameMax}{+0.187}
\newcommand{\BciFrameBestKernel}{Fidelity}
\newcommand{\BciFrameMaxP}{0.004}
\newcommand{\BciFrameAllBetter}{yes}
\newcommand{\BciFrameN}{9}
\newcommand{\FourQFrameMin}{+0.075}
\newcommand{\FourQFrameMax}{+0.106}
\newcommand{\FourQFrameBestKernel}{HS overlap}
\newcommand{\FourQFrameMaxP}{0.004}
\newcommand{\FourQFrameAllBetter}{no}
\newcommand{\FourQFrameN}{30}
\newcommand{\FourQTwinName}{control/riemann-kernel-SVM}
\newcommand{\FourQTwinAcc}{0.625}
\newcommand{\FourQBestKernel}{quantum/Bures-RBF-ref-SVM}
\newcommand{\FourQBestAcc}{0.626}
\newcommand{\FourQTwinDeltaMin}{-0.0114}
\newcommand{\FourQTwinDeltaMax}{+0.0005}
\newcommand{\FourQTwinMinP}{\ensuremath{{}={}}0.253}
\newcommand{\FourQNSubjects}{30}
\newcommand{\FourQTwinTopKernel}{quantum/Bures-RBF-ref-SVM}
\newcommand{\FourQTwinTopDelta}{+0.0005}
\newcommand{\FourQTwinTopP}{\ensuremath{{}={}}0.650}
\newcommand{\FourQTwinTopBetter}{14/30}
\newcommand{\FourQTwinRestDeltaMin}{-0.0114}
\newcommand{\FourQTwinRestDeltaMax}{-0.0015}
\newcommand{\FourQTwinRestMinP}{\ensuremath{{}={}}0.253}
\newcommand{\FourQBestClassical}{classical/TS+LR}
\newcommand{\FourQBestClassicalAcc}{0.606}
\newcommand{\FourQHeadRefDelta}{+0.0200}
\newcommand{\FourQHeadRefP}{\ensuremath{{}={}}0.016}
\newcommand{\FourQHeadRefBetter}{21/30}
\newcommand{\FourQTwinHeadDelta}{+0.0195}
\newcommand{\FourQTwinHeadP}{\ensuremath{{}={}}0.017}
\newcommand{\FourQTwinHeadBetter}{17/30}
\newcommand{\ChoFrameMin}{+0.012}
\newcommand{\ChoFrameMax}{+0.063}
\newcommand{\ChoFrameBestKernel}{HS overlap}
\newcommand{\ChoFrameMaxP}{0.009}
\newcommand{\ChoFrameAllBetter}{no}
\newcommand{\ChoFrameN}{52}
\newcommand{\ChoTwinName}{control/riemann-kernel-SVM}
\newcommand{\ChoTwinAcc}{0.638}
\newcommand{\ChoBestKernel}{quantum/Fidelity-ref-SVM}
\newcommand{\ChoBestAcc}{0.648}
\newcommand{\ChoTwinDeltaMin}{-0.0008}
\newcommand{\ChoTwinDeltaMax}{+0.0100}
\newcommand{\ChoTwinMinP}{\ensuremath{{}={}}0.001}
\newcommand{\ChoNSubjects}{52}
\newcommand{\ChoTwinTopKernel}{quantum/Fidelity-ref-SVM}
\newcommand{\ChoTwinTopDelta}{+0.0100}
\newcommand{\ChoTwinTopP}{\ensuremath{{}={}}0.001}
\newcommand{\ChoTwinTopBetter}{39/52}
\newcommand{\ChoTwinRestDeltaMin}{-0.0008}
\newcommand{\ChoTwinRestDeltaMax}{+0.0041}
\newcommand{\ChoTwinRestMinP}{\ensuremath{{}={}}0.128}
\newcommand{\ChoBestClassical}{classical/TS+LR}
\newcommand{\ChoBestClassicalAcc}{0.644}
\newcommand{\ChoHeadRefDelta}{+0.0039}
\newcommand{\ChoHeadRefP}{\ensuremath{{}={}}0.174}
\newcommand{\ChoHeadRefBetter}{31/52}
\newcommand{\ChoTwinHeadDelta}{-0.0061}
\newcommand{\ChoTwinHeadP}{\ensuremath{{}={}}0.010}
\newcommand{\ChoTwinHeadBetter}{18/52}
\newcommand{\ChoNTrialsMin}{200}
\newcommand{\ChoNTrialsMax}{240}
\newcommand{\ChoNTrialsMaxCount}{3}
\newcommand{\ChoHeadSensorDelta}{+0.0152}
\newcommand{\ChoHeadSensorP}{\ensuremath{{}<{}}0.001}
\newcommand{\ChoSensorBestQuantum}{quantum/QRE-RBF-SVM}
\newcommand{\SeedRbfWorst}{-0.0284}
\newcommand{\SeedRbfMinP}{\ensuremath{{}<{}}0.001}
\newcommand{\SeedOverlapDeltaMin}{-0.0079}
\newcommand{\SeedOverlapDeltaMax}{+0.0081}
\newcommand{\SeedOverlapMinP}{\ensuremath{{}={}}0.194}
\newcommand{\SeedSigCount}{5}
\newcommand{\SeedSigN}{15}
\newcommand{\SeedSigFavourQuantum}{0}
\newcommand{\SeedSigFavourTwin}{5}
\newcommand{\SeedN}{3}
\newcommand{\SeedFrameMin}{+0.032}
\newcommand{\SeedFrameMax}{+0.107}
\newcommand{\SeedTwinDeltaMin}{-0.0284}
\newcommand{\SeedTwinDeltaMax}{+0.0081}
\newcommand{\SeedTwinMinP}{\ensuremath{{}<{}}0.001}
\newcommand{\SeedEquivPass}{4}
\newcommand{\SeedEquivN}{15}
\newcommand{\SeedWorstBound}{0.042}
\newcommand{\SeedHeadRefMin}{+0.0183}
\newcommand{\SeedHeadRefMax}{+0.0215}
\newcommand{\SeedHeadReversals}{3}
\newcommand{\PhysCostTwinSec}{3.2}
\newcommand{\PhysCostSlowName}{quantum/Bures-RBF-ref-SVM}
\newcommand{\PhysCostSlowSec}{12.1}
\newcommand{\PhysCostRatio}{3.8}
\newcommand{\BciCostTwinSec}{15.7}
\newcommand{\BciCostSlowName}{quantum/Bures-RBF-ref-SVM}
\newcommand{\BciCostSlowSec}{117.7}
\newcommand{\BciCostRatio}{7.5}
\newcommand{\FbFrameMin}{+0.124}
\newcommand{\FbFrameMax}{+0.152}
\newcommand{\FbFrameBestKernel}{Fidelity}
\newcommand{\FbFrameMaxP}{\ensuremath{<}0.001}
\newcommand{\FbFrameAllBetter}{no}
\newcommand{\FbFrameN}{30}
\newcommand{\FbcspDelta}{+0.0212}
\newcommand{\FbcspP}{\ensuremath{{}={}}0.125}
\newcommand{\FbcspBestAcc}{0.648}
\newcommand{\FbcspAcc}{0.627}
\newcommand{\TransferN}{30}
\newcommand{\TransferSpread}{0.0141}
\newcommand{\TransferBest}{quantum/HS-overlap}
\newcommand{\TransferBestAcc}{0.647}
\newcommand{\TransferFloor}{1.9\times10^{-9}}
\newcommand{\CsN}{9}
\newcommand{\CsFloor}{3.9\times10^{-3}}
\newcommand{\CsTrialsPerSession}{144}
\newcommand{\CsFrameQMin}{+0.169}
\newcommand{\CsFrameQMax}{+0.217}
\newcommand{\CsFrameQMaxP}{0.008}
\newcommand{\CsFrameQAllBetter}{no}
\newcommand{\CsSensorQMin}{0.542}
\newcommand{\CsSensorQMax}{0.578}
\newcommand{\CsRefQMin}{0.746}
\newcommand{\CsRefQMax}{0.760}
\newcommand{\CsBestKernel}{quantum/Bures-RBF}
\newcommand{\CsBestAcc}{0.760}
\newcommand{\CsFrameClMin}{+0.022}
\newcommand{\CsFrameClMax}{+0.031}
\newcommand{\CsSensorClMin}{0.701}
\newcommand{\CsSensorClMax}{0.717}
\newcommand{\CsBestClassical}{classical/TS+LR}
\newcommand{\CsBestClassicalAcc}{0.738}
\newcommand{\CsTwinName}{control/riemann-kernel-SVM}
\newcommand{\CsTwinAcc}{0.755}
\newcommand{\CsTwinDeltaMin}{-0.0085}
\newcommand{\CsTwinDeltaMax}{+0.0054}
\newcommand{\CsTwinMinP}{\ensuremath{{}={}}0.500}
\newcommand{\CsTwinAbsMax}{0.0085}
\newcommand{\CsSpread}{0.0139}
\newcommand{\SweepQMin}{3}
\newcommand{\SweepQMax}{6}
\newcommand{\SweepN}{14}
\newcommand{\SweepRefRatioMin}{1.6}
\newcommand{\SweepRefRatioMax}{2.4}
\newcommand{\SweepSensorRatioMin}{4.1}
\newcommand{\SweepSensorRatioMax}{7.7}
\newcommand{\SweepGainLoMin}{4.7}
\newcommand{\SweepGainLoMax}{9.5}
\newcommand{\SweepGainHiMin}{1.9}
\newcommand{\SweepGainHiMax}{2.6}
\newcommand{\SweepRefAboveBase}{every}
\newcommand{\SweepRefMonotone}{not every}
\newcommand{\SweepRefDipMax}{8}
\newcommand{\GramGainMin}{4.7}
\newcommand{\GramGainMax}{9.5}
\newcommand{\GramNSubjects}{14}
\newcommand{\GramSensorMean}{0.981}
\newcommand{\GramSensorStd}{0.022}
\newcommand{\ShotCrossover}{10^{4}}
\newcommand{\FrameEffectMin}{+0.012}
\newcommand{\FrameEffectMax}{+0.217}


%% ------------------------------ equivalence macros

\newcommand{\EquivMargin}{0.02}
\newcommand{\EquivN}{35}
\newcommand{\EquivPass}{20}
\newcommand{\EquivWorstBound}{0.032}
\newcommand{\EquivFailQuantum}{10}
\newcommand{\EquivFailClassical}{5}
\newcommand{\EquivNZeroExcluded}{3}
\newcommand{\EquivNSettings}{seven}
\newcommand{\EquivNDatasets}{three}
\newcommand{\EquivNRegisters}{three}
\newcommand{\EquivBciPass}{5}
\newcommand{\EquivBciN}{5}
\newcommand{\EquivBciWorst}{0.0109}
\newcommand{\EquivCsPass}{2}
\newcommand{\EquivCsN}{5}
\newcommand{\EquivCsWorst}{0.0281}
\newcommand{\EquivTrPass}{3}
\newcommand{\EquivTrN}{5}
\newcommand{\EquivTrWorst}{0.0235}
\newcommand{\EquivFourQPass}{4}
\newcommand{\EquivFourQN}{5}
\newcommand{\EquivFourQWorst}{0.0263}
\newcommand{\EquivChoPass}{5}
\newcommand{\EquivChoN}{5}
\newcommand{\EquivChoWorst}{0.0153}
\newcommand{\EquivZeroExclName}{Fidelity (PhysioNet, 5 qubits)}
\newcommand{\EquivZeroExclDelta}{+0.0180}
\newcommand{\EquivZeroExclList}{Fidelity on PhysioNet, 5 qubits at $+0.0180$, HS overlap on Transfer (LOSO) at $+0.0119$ and Fidelity on Cho2017 at $+0.0100$}
\newcommand{\EquivZeroExclAllQuantum}{yes}


\begin{document}

%% iopjournal prints the article type under a "Journal Name" placeholder that
%% IOP replace in production. JNE article types: Paper, Note, Topical Review,
%% Comment, Reply.
\articletype{Paper}

%% \title takes ONE argument here; iopart's optional short-title form is not
%% supported by this class.
\title{Quantum Kernels for EEG: An Invariance Mismatch Masquerading as
Advantage}

%% AUTHOR ORDER IS NOT FINAL -- CONFIRM BEFORE SUBMITTING.
%% Order below is the usual convention (student first, supervisor last).
%% All three names were read from the ORCID public records rather than typed:
%%   0009-0009-2423-1835 -> "Md. Imtiaj Alam Sajin"
%%   0009-0009-7116-7151 -> "Samiul Islam Saif"
%%   0000-0003-3741-6651 -> "Muhammad Hasibur Rashid Chayon"
%% Chayon's ORCID lists two employers (AIUB and IUBAT); AIUB is used here as
%% the institution where the work was supervised. Confirm.
\author{Md. Imtiaj Alam Sajin$^{1,*}$\orcid{0009-0009-2423-1835},
Samiul Islam Saif$^{1}$\orcid{0009-0009-7116-7151} and
Muhammad Hasibur Rashid Chayon$^{1}$\orcid{0000-0003-3741-6651}}

\affil{$^1$Department of Computer Science, American International
University-Bangladesh (AIUB), Dhaka, Bangladesh}

\affil{$^*$Author to whom any correspondence should be addressed.}

%% A student address expires; the gmail is the durable contact and is listed
%% second so the correspondence route survives graduation.
\email{26-94090-2@student.aiub.edu, imtiajsajin@gmail.com (M I A Sajin);
Md.samiulislam.aiub@gmail.com (S I Saif); chayon@aiub.edu (M H R Chayon)}

\begin{abstract}
\textit{Objective.}
Quantum machine learning is often proposed for EEG decoding, usually without
the controls that show the quantum part is doing the work.
We ask whether quantum kernels beat well-tuned classical baselines, and if not,
why.

\textit{Approach.}
A trace-normalised EEG covariance matrix is a quantum density matrix, so
quantum information geometry applies directly. We benchmark Hilbert--Schmidt,
fidelity, Bures and quantum-relative-entropy kernels against tuned Riemannian,
CSP and filter-bank baselines on two motor-imagery datasets under one nested
cross-validation budget. One asymmetry stands out:
the quantum quantities are invariant only under orthogonal congruence, while
EEG's nuisances span the full congruence group the classical baselines already
handle. Borrowing recentring from Riemannian transfer learning, we refer each
state to a training-set Fr\'echet mean, prove this restores exact affine
invariance, and add a twin differing only in metric.

\textit{Main results.}
In the sensor frame the classical baselines won on both datasets. Recentring
raised Gram variance several-fold and improved every kernel, by up to
\PhysFrameMax{} on PhysioNet and \BciFrameMax{} on IV-2a, the latter in every
subject, reversing the headline comparison from $\HeadSensorDelta$ for the best
classical pipeline to $\HeadRefDelta$ against it. The twin control shows this
is not quantum. Across \EquivNSettings{} settings (\EquivNDatasets{} datasets,
\EquivNRegisters{} register sizes, within-subject, cross-subject and
cross-session evaluation), two one-sided
tests place the two families within $\pm\EquivWorstBound$ accuracy in all
\EquivN{} comparisons; under two further cross-validation partitions the
bandwidth-tuned quantum kernels fall behind the twin, never ahead. Parity held
against filter-bank CSP and in cross-session transfer, where non-congruence
shift is largest.

\textit{Significance.}
We find no quantum advantage, and trace it to an invariance mismatch rather
than kernel concentration: supply the missing invariance and the quantum
kernels match their classical counterparts. Recentring is a precondition for a
fair comparison; omitting it is a concrete route to a false positive.
\end{abstract}

\keywords{quantum machine learning, brain--computer interface,
electroencephalography, quantum kernels, Riemannian geometry, motor imagery,
benchmarking}

%% =====================================================================
\section{Introduction}
%% =====================================================================

Decoding motor imagery from EEG is a core problem in brain--computer interface
(BCI) research, and the field has converged on a small set of strong classical
methods: common spatial patterns (CSP) \cite{ramoser2000} and, more recently,
classifiers built on the Riemannian geometry of spatial covariance matrices
\cite{barachant2012,congedo2017}. Large-scale benchmarking finds that
covariance tangent-space projection and CSP consistently define the strongest
methodological families across public motor-imagery datasets; Lotte \textit{et al}
\cite{lotte2018} and Yger \textit{et al} \cite{yger2017} review the field and the
Riemannian branch of it respectively.

In parallel, a growing literature applies quantum machine learning (QML) to
EEG, typically reporting improvements over a classical comparator
\cite{qeegnet2024,eswa2024}. These reports share a recurring methodological
pattern that makes them difficult to interpret. Classical baselines are often
weak or untuned: plain support vector machines on band power rather than the
Riemannian or filter-bank CSP methods that define the state of the art.
Hyperparameters are frequently optimised for the quantum model alone. Most
critically, two controls are almost universally absent: an \emph{entanglement
ablation}, which tests whether the quantum resource contributes anything at
all, and a \emph{dimension-matched classical comparator}, without which an
apparent quantum effect is confounded with the dimensionality reduction applied
before encoding.

Quantum kernels are the natural object of study here. Supervised quantum
models with fixed encodings are kernel methods \cite{schuld2019,
schuld2021kernel}, so a kernel comparison captures the whole model class
rather than one architecture, and it sidesteps the trainability problems
that afflict variational circuits \cite{mcclean2018,cerezo2021}. What is
known about their usefulness is mixed. There is a constructed problem with
a proven exponential separation \cite{liu2021}, but on natural data the
advantage is bounded by how much the quantum kernel's inductive bias
happens to match the target \cite{kubler2021,huang2021}, and general
surveys \cite{biamonte2017} have not translated into demonstrated wins on
biological signals.

The wider QML literature gives strong reasons for caution. Bowles \textit{et al}
\cite{bowles2024} benchmarked twelve quantum models across 160 datasets and
found that out-of-the-box classical models generally outperformed them; more
pointedly, removing entanglement often left performance equal or better.
Thanasilp \textit{et al} \cite{thanasilp2024} proved that quantum kernel values
concentrate exponentially in qubit count toward a fixed value, collapsing the
model to a trivial one. Neither result has been systematically tested on EEG.

This paper makes four contributions.

First, we identify a \emph{principled} quantum encoding for EEG. An EEG trial's
spatial covariance matrix is symmetric positive definite; normalised to unit
trace it is a bona fide quantum density matrix (figure~\ref{fig:eeg}). This is
an identity, not an analogy, and it makes quantum information geometry (Uhlmann
fidelity, the Bures metric, von Neumann entropy) applicable to EEG natively,
with no lossy squeezing of features into rotation angles. The Hilbert--Schmidt overlap
$\mathrm{tr}(\rho\sigma)$ is moreover exactly the quantity a SWAP test
estimates, so the construction maps onto a hardware-native primitive. Related
tooling exists \cite{pyriemannqiskit2023}, but its quantum step operates on
flattened tangent-space vectors, discarding the density-matrix structure.

Second, we report a controlled benchmark of \NPipelinesExt{} pipelines on
\NSubjects{} subjects with nested cross-validation, matched tuning budgets,
both of the missing controls, and paired subject-wise statistics with
multiple-comparison correction, replicated on IV-2a and extended to
filter-bank baselines, cross-subject and cross-session transfer, and a
register sweep to six qubits in both frames.

Third, and centrally, we show that the comparison such benchmarks perform, 
ours included, in its first form, is confounded by the coordinate frame. The
quantum quantities above are invariant under orthogonal congruence only, while
EEG's nuisance transformations (electrode gain, reference choice, source
mixing, and the change from one subject to the next) generate the full
congruence group, under which every strong classical baseline \emph{is}
invariant. We prove that referring each state to a Fr\'echet-mean reference
state restores exact affine invariance while remaining a legitimate quantum
operation, and we show that this single change relieves concentration by
\GramGainMin--\GramGainMax$\times$ and reverses the sign of the
classical-versus-quantum comparison on both datasets.

Fourth, we show that the reversal is nonetheless \emph{not} evidence of quantum
advantage, using a control we believe should be standard in this literature: a
classical SPD-manifold kernel placed in the same classifier, on the same data,
in the same frame, differing from the quantum kernels only in the metric. No
quantum kernel exceeds that twin in any of \EquivNSettings{} settings,
including cross-session transfer, the regime in which the shift our invariance
argument does not cover is proportionally largest. The
practical implication is uncomfortable but concrete: a quantum kernel compared
against a classical baseline in a different frame will show an advantage of
roughly the size we report, produced entirely by the frame.

\begin{figure}[htbp]
\centering
\includegraphics[width=\textwidth]{fig0_eeg_to_state.pdf}
\caption{\label{fig:eeg}From scalp EEG to a quantum state, computed from real
recordings rather than simulated. (a) Band-passed sensorimotor traces for one
left-hand and one right-hand trial. (b) Scalp topography of mu/beta power,
right minus left: power falls over the hemisphere contralateral to the imagined
hand, which is the physiological effect every decoder in this study exploits.
(c) The spatial covariance of a single trial over the eight decoding channels.
(d) The same matrix divided by its trace. It is now a density matrix, and
(e) its eigenvalues are non-negative and sum to one, so they are a population
distribution. The step from (c) to (d) is the whole encoding: no dimensionality
reduction, no rotation-angle embedding.}
\end{figure}

We state our scope plainly. At the channel counts considered here (8--64,
i.e.\ 3--6 qubits) every kernel we evaluate is classically computable in
$O(n_{\mathrm{ch}}^3)$. We make no claim of quantum speedup; the object of
study is whether quantum \emph{geometry} improves decoding. Every quantum
kernel is evaluated by exact state-vector simulation: there is no gate error,
no decoherence and no readout bias, and section~\ref{sec:shots} adds only sampling
noise. Hardware would make the quantum side worse, never better, so the
comparison throughout is maximally generous to it and a negative result here
is not weakened by the idealisation.

The contribution we intend for neural engineering is practical rather than
quantum. Covariance decoders are the field's working standard, and the two
devices this paper turns on the quantum kernels, recentring to a label-free
reference state and a comparator that differs in metric alone, apply to any
new covariance kernel a practitioner might be offered, quantum or otherwise.
The invariance proposition says \emph{why} recentring is a precondition for
comparing geometries at all, where the Riemannian transfer-learning literature
\cite{zanini2018,rodrigues2019} has so far justified it empirically; and the
negative result is stated with enough control that it can be built on rather
than merely cited.

%% =====================================================================
\section{Methods}
%% =====================================================================

\subsection{Data and preprocessing}

We use the PhysioNet EEG Motor Movement/Imagery Database (EEGMMIDB)
\cite{schalk2004,goldberger2000}, recorded with the BCI2000 system at 160~Hz
over 64 channels. We analyse the standard binary protocol: imagined left-hand
versus right-hand movement (runs 4, 8 and 12), giving \NTrials{} trials per
subject. Signals are band-pass filtered to 8--30~Hz, spanning the $\mu$ and
$\beta$ rhythms in which motor-imagery event-related desynchronisation appears,
and epoched from 0.5 to 3.5~s after cue onset to exclude the cue-evoked
transient. Epochs are resampled to 128~Hz.

The primary analysis uses eight sensorimotor channels (FC3, FCz, FC4, C3, Cz,
C4, CP3, CP4), a register of exactly three qubits under the amplitude encoding
described below. Subjects 88, 89, 92, 100 and 104 are excluded for documented
annotation and sampling-rate defects. We analyse the first \NSubjects{} eligible
subjects.

Spatial covariance matrices are estimated per trial with the oracle
approximating shrinkage (OAS) estimator, which is well conditioned at the
trial counts available here.

\subsection{Second dataset: BCI Competition IV-2a}

A single dataset cannot separate a property of the methods from a property of
the data. PhysioNet gives breadth but only \NTrials{} trials per subject, so a
negative result there invites the objection that the quantum models were
starved. We therefore repeat the entire benchmark on BCI Competition IV-2a
(\BciNSubjects{} subjects, \BciNTrials{} trials each, \BciNFoldScores{} fold
scores), accessed through MOABB \cite{moabb2018}. The two datasets sit at
opposite corners of the design space: many subjects with little data each
versus few subjects with a great deal.

IV-2a is also the dataset the quantum-EEG literature reports on
\cite{qeegnet2024}, so results here are directly comparable to published
claims. The same eight sensorimotor channels are used, all of which exist in
the IV-2a montage, so the register size is identical at three qubits.
Band-pass, resampling, pipelines, hyperparameter grids and cross-validation
protocol are unchanged; only the data differs. Each dataset keeps its own
standard epoch window, since each competition defines its own cue timing, and
the comparison of interest is whether the \emph{ranking} of pipelines moves rather
than whether absolute accuracies coincide.

\subsection{EEG covariance matrices as quantum states}

Let $C_i \in \mathbb{R}^{n_{\mathrm{ch}} \times n_{\mathrm{ch}}}$ be the
spatial covariance of trial $i$. Since $C_i$ is symmetric positive definite,
\begin{equation}
\rho_i \;=\; \frac{C_i}{\mathrm{tr}\,C_i}
\label{eq:density}
\end{equation}
satisfies $\rho_i \succeq 0$ and $\mathrm{tr}\,\rho_i = 1$: it is a density
matrix on an $n_{\mathrm{ch}}$-dimensional Hilbert space, that is, on
$\log_2 n_{\mathrm{ch}}$ qubits. Sweeping the channel count over powers of two
therefore sweeps the register size with no change of method, a property we
exploit in section~\ref{sec:concentration}.

\subsection{Density-matrix kernels}

Equation~(\ref{eq:density}) makes the following quantum-information quantities
directly available as kernels between trials.

The \emph{Hilbert--Schmidt overlap}
\begin{equation}
K_{\mathrm{HS}}(\rho,\sigma) = \mathrm{tr}(\rho\sigma)
\end{equation}
is an inner product on Hermitian matrices, so its Gram matrix is positive
semi-definite by construction and requires no repair. It is precisely the
observable estimated by a SWAP test.

The \emph{Uhlmann fidelity} \cite{uhlmann1976}
\begin{equation}
F(\rho,\sigma) = \left(\mathrm{tr}\sqrt{\sqrt{\rho}\,\sigma\sqrt{\rho}}\right)^{2}
\end{equation}
is the quantum generalisation of the Bhattacharyya coefficient. We evaluate it
via the eigenvalues of the symmetrised product, which is numerically stable for
the matrix sizes involved.

Both quantities proved severely concentrated on EEG (section~\ref{sec:concentration}),
so we additionally define bandwidth-parameterised kernels that retain the
quantum geometry while restoring a length scale, exactly as a radial basis
function (RBF) kernel does for the Euclidean distance:
%% iopart.cls is not compatible with amsmath's align; eqnarray is the
%% class-supported multi-line display.
\begin{eqnarray}
K_{\mathrm{HS\mbox{-}RBF}}(\rho,\sigma) &=& \exp\left(-\gamma\,
   \mathrm{tr}\left[(\rho-\sigma)^2\right]\right), \\
K_{\mathrm{Bures}}(\rho,\sigma) &=& \exp\left(-\gamma\,
   d_{B}^{2}(\rho,\sigma)\right), \qquad
d_{B}^{2} = 2\left(1-\sqrt{F(\rho,\sigma)}\right).
\end{eqnarray}
Here $d_B$ is the Bures distance \cite{bhatia2019}, the geodesic distance
of the Bures metric on
state space, which for unit-trace matrices coincides with the
Bures--Wasserstein distance occasionally used in the Riemannian EEG literature
as a distance but, to our knowledge, never as a kernel. The bandwidth
$\gamma$ is a bandwidth, and bandwidth is known to control quantum-kernel
expressivity and to be decisive for whether such kernels generalise at all
\cite{shaydulin2022}. It is set by the median heuristic computed
\emph{on the training split
only}, with a multiplier selected in the inner cross-validation loop.

\subsection{Reference states, and the invariance the sensor frame lacks}
\label{sec:reference}

Equation~(\ref{eq:density}) expresses each state in the frame the electrodes
happen to define. That frame is not neutral.

The nuisance transformations of EEG act on the spatial covariance by
congruence, $C \mapsto ACA^{\mathsf{T}}$ with $A$ invertible: per-electrode
gain and impedance, the choice of reference electrode, linear source mixing
and volume conduction, and the change from one subject or session to the next
all take this form. The affine-invariant Riemannian metric is invariant under
the whole congruence group, which is the structural reason tangent-space
decoding is the classical state of the art and why it transfers.

The quantum quantities are not. $\mathrm{tr}(\rho\sigma)$, the Uhlmann
fidelity, the Bures distance and the quantum relative entropy are invariant
only under the \emph{orthogonal} subgroup $A \in O(n_{\mathrm{ch}})$. This is
an asymmetry in the comparison itself, not merely a property of the data:
pyRiemann's tangent-space estimator fits a reference mean $M$ and maps
$C \mapsto \log(M^{-1/2} C M^{-1/2})$, and both the minimum-distance-to-mean
classifier and CSP are congruence-invariant by construction. Every strong
classical baseline in our suite therefore possesses an invariance that none of
the sensor-frame quantum kernels possess.

The repair is not new, and we want to be plain about that. Whitening
covariances by their Fr\'echet mean is \emph{Riemannian recentring}, introduced
by Zanini \textit{et al} \cite{zanini2018} to align sessions and subjects and extended
by Rodrigues \textit{et al} \cite{rodrigues2019} into Riemannian Procrustes Analysis;
He and Wu \cite{he2020} give a Euclidean counterpart. In the classical setting
it is understood as a domain-adaptation step, applied because it helps in
transfer, and it is optional: the affine-invariant methods it is applied to are
already congruence-invariant without it.

Our point is that for quantum kernels the same operation plays a different and
stronger role. It is not an optional aid to transfer but the step that supplies
an invariance the kernels structurally lack, and which the baselines they are
being compared against already have. We therefore refer each state to a
\emph{reference state} $M$, the Fr\'echet mean of the \emph{training}
covariances, via the whitening $W = M^{-1/2}$:
\begin{equation}
\tilde{\rho}_i \;=\; \frac{W C_i W}{\mathrm{tr}\,(W C_i W)}.
\label{eq:refstate}
\end{equation}

%% iopart.cls defines no theorem environments and loading amsthm alongside it
%% is not worth the risk; the statement is set in the class's own \indented.
\textbf{Proposition.}\ \emph{Under a common congruence
$C_i \mapsto A C_i A^{\mathsf{T}}$ with $A$ invertible, the kernels
$\mathrm{tr}(\tilde\rho\tilde\sigma)$, $F(\tilde\rho,\tilde\sigma)$,
$d_B(\tilde\rho,\tilde\sigma)$ and $J(\tilde\rho,\tilde\sigma)$ are
invariant.}

\noindent\emph{Proof.} The affine-invariant mean is equivariant,
$M \mapsto AMA^{\mathsf{T}}$. Put $B = (AMA^{\mathsf{T}})^{-1/2}A$. Then
$BMB^{\mathsf{T}} = I = WMW^{\mathsf{T}}$, so $U := BM^{1/2}$ satisfies
$UU^{\mathsf{T}} = I$ and $B = UW$. Hence
$W C_i W \mapsto U (W C_i W) U^{\mathsf{T}}$ for one common orthogonal $U$,
and trace normalisation commutes with orthogonal conjugation. Each quantity
above is invariant under a common unitary. \hfill$\square$

In the reference frame the quantum kernels are thus \emph{exactly}
affine-invariant, sharing the invariance group of the affine-invariant
Riemannian metric. Figure~\ref{fig:reference} sets out the geometry and
verifies the invariance numerically. Two consequences follow, and we test both. First, the
comparison against the classical baselines becomes like for like. Second, and
less obviously, the quantum kernels can now at best \emph{match} the
Riemannian metric on any nuisance that is a congruence; parity, not
superiority, is the ceiling the algebra permits.

To state the contribution precisely: the alignment is borrowed, the reason for
needing it is not. What we add is the observation that quantum kernels lack an
invariance the classical baselines already possess, a proof that recentring
supplies exactly that invariance for the quantum quantities, and the control in
section~\ref{sec:refresults} showing that once it is supplied nothing quantum remains
to distinguish them.

Equation~(\ref{eq:refstate}) is not classical preprocessing smuggled into a
quantum method. The map $\rho \mapsto W\rho W/\mathrm{tr}(W\rho W)$ is a
filtering (L\"uders) operation with Kraus operator $W$ followed by
renormalisation, so it is state preparation and the SWAP-test implementation
path is unchanged. $M$ is estimated on training data only.

\begin{figure}[htbp]
\centering
\includegraphics[width=\textwidth]{fig6_reference.pdf}
\caption{\label{fig:reference}(a) EEG covariances cluster tightly about their
Fr\'echet mean $M$, and a nuisance congruence moves the whole cluster.
Whitening by $M$ spreads the states about the identity and reduces the nuisance
to a common rotation, to which every quantum invariant is blind. (b) The
proposition verified numerically: reference-frame residuals sit at machine
precision, sensor-frame ones are $O(10^{-1})$. (c) Off-diagonal Gram variance
rises by \GramGainMin--\GramGainMax$\times$ in the reference frame. (d) The
sensor-frame overlaps saturate near unity in every subject
($n=\GramNSubjects$).}
\end{figure}

We also add the symmetrised quantum relative entropy
\begin{equation}
J(\rho,\sigma) = \frac{1}{2}\,\mathrm{tr}\!\left[(\rho-\sigma)
   (\log\rho - \log\sigma)\right],
\end{equation}
used as $\exp(-\gamma J)$. It is distinct from the classical
Kullback--Leibler divergence between the Gaussians the covariances
parameterise, which pyRiemann supplies as a distance and which involves
$\mathrm{tr}(\sigma^{-1}\rho)$ and log-determinants rather than matrix
logarithms.

\subsection{Circuit embedding kernels and the entanglement ablation}

For comparison with the prevailing approach, we also evaluate
parameterised-circuit embedding kernels
$K(x,z) = |\langle \phi(x) | \phi(z) \rangle|^{2}$ (figure~\ref{fig:circuits}),
where $\phi$ is an
instantaneous-quantum-polynomial (IQP) style feature map of the
Havlíček \textit{et al} family \cite{havlicek2019}, built from Hadamard and $R_Z$
rotations with data-dependent $ZZ$ couplings. Inputs are Riemannian
tangent-space features reduced by PCA to the number of qubits and scaled to a
rotation-angle range; the angle scale is itself a tuned hyperparameter. A
variant with fixed CNOT entanglers is also included.

The \textbf{entanglement ablation} is the decisive control: an otherwise
identical circuit with every entangling gate deleted. The resulting kernel
factorises into single-qubit kernels and is classically trivial. If it matches
or exceeds the entangled version, quantum structure is not contributing
\cite{bowles2024}.

Kernels are evaluated by preparing each embedded state once on a state-vector
simulator, which is mathematically identical to running the adjoint kernel-%
estimation circuit at infinite shots but costs $O(n\,2^{q})$ rather than
$O(n^{2})$ circuit executions.

\begin{figure}[htbp]
\centering
\includegraphics[width=\textwidth]{fig5_circuits.pdf}
\caption{\label{fig:circuits}The circuit feature maps, at three qubits with one
layer shown. Rotation angles carry the PCA-reduced tangent-space features and
the kernel is the state overlap $|\langle\phi(x)|\phi(z)\rangle|^{2}$. The IQP
map uses data-dependent $ZZ$ couplings; the ring map uses fixed CNOTs. The
third panel is the entanglement ablation: the identical circuit with every
entangling gate deleted, which factorises the kernel into single-qubit terms
and isolates what entanglement contributes. The figure is rendered by Qiskit
from a circuit that is checked against the executed PennyLane circuit before
drawing: the two agree on full Gram matrices to $<10^{-15}$, so the diagram
cannot drift from the code it documents.}
\end{figure}

\subsection{Classical baselines and controls}

Five classical baselines represent the genuine state of the art rather than a
strawman: log-variance with shrinkage LDA; CSP with shrinkage LDA; minimum
distance to Riemannian mean (MDM); Riemannian tangent space with logistic
regression; and tangent space with an RBF SVM.

Four controls isolate specific confounds: the entanglement ablation described
above; \emph{dimension-matched} linear and RBF SVMs operating on exactly the
same PCA features the circuit kernels receive; and a log-Euclidean tangent-space
classifier as a first classical geometric comparator for the density-matrix
kernels.

The control the argument ultimately rests on is stricter than that, and we
define it here because everything in section~\ref{sec:refresults} depends on
it. A \emph{metric-matched twin} is a classical kernel on the SPD manifold used
exactly as the density-matrix kernels are used: the same covariances, the same
reference-state whitening (\ref{eq:refstate}), a precomputed Gram matrix in
the same support vector machine, and the same hyperparameter grid. Only the
metric differs. We use the two SPD kernels pyRiemann provides
\cite{barachant2013}: the affine-invariant Riemannian kernel,
$\exp(-\gamma\,\delta_R^2)$ with $\delta_R$ the affine-invariant distance, and
the log-Euclidean kernel. Whichever of the two scores higher in a given setting
is the comparator, which is the conservative choice. Because the twin shares
every element of the quantum pipeline except the metric, any difference
between them is attributable to the metric alone; because it is
congruence-invariant by construction, it is the right comparator for the
reference-frame kernels specifically.

Three further suites reuse this design. A \emph{filter-bank} suite splits the
8--30 Hz band into four sub-bands (8--12, 12--16, 16--22 and 22--30 Hz) and
stacks them, so that one covariance spans $4\times 8=32$ dimensions --- a
single density matrix on five qubits --- and adds filter-bank CSP
\cite{ang2008} and filter-bank tangent-space baselines alongside
filter-bank versions of every kernel and twin. A \emph{cross-subject} suite
pools the trials of all but one subject, trains on the pool, and tests on the
held-out subject; its hyperparameters are chosen by subject-grouped inner
cross-validation on the training subjects, so no held-out subject influences
any choice. A \emph{cross-session} suite uses the two recording sessions of
IV-2a: for each subject, and in both directions, it trains on one session and
tests on the other, choosing hyperparameters by stratified inner
cross-validation within the training session alone. In its reference frame
each session is whitened by its own Fr\'echet mean, which uses no labels and
is, for the test session, exactly the unsupervised adaptation a deployed
decoder can perform on the unlabelled data it receives before decoding
anything. All three suites are described where their results appear.

\subsection{Evaluation protocol and statistics}

Each subject is evaluated independently. The outer loop is a repeated
stratified $k$-fold ($5$ folds $\times\,3$ repeats) estimating generalisation;
the inner loop is a $4$-fold grid search selecting hyperparameters. Crucially,
\emph{every} pipeline, including all classical baselines, receives the same
inner search budget. All fitted transformations (tangent-space reference mean,
scalers, PCA, and the kernel bandwidth) are refit within each training split,
so no information leaks across folds.

Across subjects we follow the standard recommendation for comparing
classifiers over multiple datasets \cite{demsar2006} and use the Wilcoxon
signed-rank test on paired per-subject
accuracies, with Holm correction \cite{holm1979} across the family of
comparisons, and report
the paired effect size $d_z$. Per-subject optimality varies substantially in
motor-imagery decoding, so paired non-parametric testing (rather than
comparison of grand means) is the appropriate inference.

Two points about the correction structure, since several families appear
below. Holm correction is applied \emph{within} each family, and the families
are the pre-specified ones: the suite-wide comparison against the reference
pipeline, the pre-specified contrasts of table~\ref{tab:key}, and the
metric-matched twin comparisons. These were fixed by the design of the study
--- one family per question asked --- and not selected after inspecting
results; no comparison reported here was chosen because it was significant. We
deliberately do not apply a single global correction across every family,
because the questions are logically distinct and a global correction would
make the null results, which are the paper's actual claims, easier rather than
harder to obtain.

The outer partition is itself a source of variation, so the PhysioNet
three-qubit suite, on which the decisive comparison is first made, is
repeated in full under two further outer-partition seeds
(table~\ref{tab:seeds}); nothing else changes between repetitions.

Where the claim is that two methods are the \emph{same}, a non-significant
difference is not sufficient evidence, so we additionally report two one-sided
tests (TOST) on the paired differences against a margin fixed in advance
(table~\ref{tab:equiv}). The margin, its justification and the achieved
equivalence bounds are given there.

\subsection{Implementation}

Analyses use MNE-Python \cite{mne2013}, pyRiemann, scikit-learn and PennyLane
\cite{pennylane2018}. Code and results are available at
\url{https://github.com/Imtiaj-Sajin/quantaEEG}.

%% =====================================================================
\section{Results}
%% =====================================================================

%% Auto-generated table floats; LaTeX places them near their references.
%% AUTO-GENERATED by paper/make_tables.py: DO NOT EDIT BY HAND.
%% Regenerate after every benchmark run:  python paper/make_tables.py


%% ---------------------------------------------------------------- Table 1
\begin{table}[htbp]
\caption{\label{tab:main}Within-subject decoding performance on PhysioNet
EEGMMIDB left- versus right-hand motor imagery. Accuracy and AUC are means over
subjects of the per-subject nested cross-validation score; SD is between
subjects. Runtime is mean wall-clock seconds per subject for the full
hyperparameter search. Rows are ordered by accuracy.}
\begin{tabular}{@{}llcccc@{}}
\hline
Pipeline & Group & Accuracy & SD & AUC & Runtime (s) \\
\hline
classical/CSP+LDA & Classical & \textbf{0.611} & 0.154 & 0.636 & 1.99 \\
classical/TS+LR & Classical & 0.595 & 0.174 & 0.625 & 2.69 \\
classical/TS+RBF-SVM & Classical & 0.590 & 0.168 & 0.605 & 6.27 \\
control/logeuclid-TS+LR & Control & 0.582 & 0.149 & 0.611 & 2.65 \\
control/PCA-matched-linear & Control & 0.582 & 0.156 & 0.588 & 2.41 \\
quantum/CNOT-kernel-SVM & Quantum & 0.579 & 0.148 & 0.593 & 9.27 \\
control/PCA-matched-RBF & Control & 0.578 & 0.153 & 0.576 & 6.64 \\
control/IQP-no-entangle & Control & 0.576 & 0.147 & 0.594 & 8.26 \\
classical/MDM & Classical & 0.575 & 0.161 & 0.594 & 0.23 \\
quantum/IQP-kernel-SVM & Quantum & 0.566 & 0.134 & 0.588 & 10.19 \\
classical/logvar+LDA & Classical & 0.559 & 0.134 & 0.585 & 0.05 \\
quantum/Bures-RBF-SVM & Quantum & 0.552 & 0.109 & 0.567 & 12.40 \\
quantum/HS-RBF-SVM & Quantum & 0.549 & 0.101 & 0.567 & 5.85 \\
quantum/HS-overlap-SVM & Quantum & 0.532 & 0.076 & 0.578 & 2.09 \\
quantum/Fidelity-SVM & Quantum & 0.523 & 0.082 & 0.557 & 4.68 \\
\hline
\end{tabular}
\end{table}


%% ---------------------------------------------------------------- Table 2
\begin{table}[htbp]
\caption{\label{tab:tests}Paired comparisons against the classical/TS+LR
reference across subjects (Wilcoxon signed-rank, two-sided). $\Delta$ is the
mean per-subject accuracy difference; $d_z$ is the paired effect size;
$p_{\mathrm{Holm}}$ is corrected across the whole family of comparisons. No comparison survives correction.}
\begin{tabular}{@{}lccccc@{}}
\hline
Pipeline & $\Delta$ accuracy & $p$ & $p_{\mathrm{Holm}}$ & $d_z$ & Better in \\
\hline
classical/CSP+LDA & $+0.0168$ & 0.156 & 1.000 & $+0.261$ & 17/30 \\
classical/TS+RBF-SVM & $-0.0047$ & 0.905 & 1.000 & $-0.105$ & 16/30 \\
control/logeuclid-TS+LR & $-0.0121$ & 0.405 & 1.000 & $-0.140$ & 13/30 \\
control/PCA-matched-linear & $-0.0126$ & 0.946 & 1.000 & $-0.165$ & 18/30 \\
quantum/CNOT-kernel-SVM & $-0.0156$ & 0.227 & 1.000 & $-0.217$ & 11/30 \\
control/PCA-matched-RBF & $-0.0170$ & 0.264 & 1.000 & $-0.242$ & 9/30 \\
control/IQP-no-entangle & $-0.0183$ & 0.316 & 1.000 & $-0.241$ & 11/30 \\
classical/MDM & $-0.0195$ & 0.198 & 1.000 & $-0.262$ & 12/30 \\
quantum/IQP-kernel-SVM & $-0.0289$ & 0.131 & 1.000 & $-0.316$ & 13/30 \\
classical/logvar+LDA & $-0.0360$ & 0.141 & 1.000 & $-0.283$ & 11/30 \\
quantum/Bures-RBF-SVM & $-0.0425$ & 0.081 & 0.978 & $-0.375$ & 12/30 \\
quantum/HS-RBF-SVM & $-0.0452$ & 0.111 & 1.000 & $-0.358$ & 13/30 \\
quantum/HS-overlap-SVM & $-0.0625$ & 0.048 & 0.628 & $-0.438$ & 12/30 \\
quantum/Fidelity-SVM & $-0.0714$ & 0.013 & 0.179 & $-0.498$ & 10/30 \\
\hline
\end{tabular}
\end{table}


%% ---------------------------------------------------------------- Table 3
\begin{table}[htbp]
\caption{\label{tab:key}Pre-specified comparisons, first-named pipeline minus
second. The entanglement ablation contrasts the IQP circuit kernel with an
otherwise identical circuit whose entangling gates are deleted; the
dimension-matched control contrasts it with a linear SVM on exactly the same
PCA features. Significant at $p<0.05$ with $n=30$: Best classical versus best quantum; the others are not.}
\begin{tabular}{@{}lcccc@{}}
\hline
Comparison & $\Delta$ accuracy & $p$ & $d_z$ & Better in \\
\hline
Best classical vs.\ best quantum & \textbf{$+0.0323$} & \textbf{0.002} & $+0.576$ & 24/30 \\
Entanglement ablation & {$+0.0106$} & {0.309} & $+0.236$ & 17/30 \\
Dimension-matched control & {$+0.0163$} & {0.173} & $+0.275$ & 18/30 \\
\hline
\end{tabular}
\end{table}


%% ---------------------------------------------------------------- Table 5
\begin{table}[htbp]
\caption{\label{tab:bci}Replication on BCI Competition IV-2a (9 subjects, 288 trials each), with the PhysioNet result repeated for
comparison. Protocol, channels, tuning budget and pipelines are
identical; only the data differs. Quantum kernels occupy the four
lowest positions on both datasets.}

\begin{tabular}{@{}lccc@{}}
\hline
Pipeline & Acc (IV-2a) & AUC (IV-2a) & Acc (Phys.) \\
\hline
classical/TS+LR & \textbf{0.763} & 0.822 & 0.595 \\
classical/TS+RBF-SVM & 0.755 & 0.811 & 0.590 \\
classical/CSP+LDA & 0.751 & 0.808 & 0.611 \\
control/logeuclid-TS+LR & 0.751 & 0.810 & 0.582 \\
classical/MDM & 0.723 & 0.777 & 0.575 \\
classical/logvar+LDA & 0.697 & 0.748 & 0.559 \\
control/PCA-matched-linear & 0.690 & 0.735 & 0.582 \\
control/PCA-matched-RBF & 0.687 & 0.729 & 0.578 \\
quantum/CNOT-kernel-SVM & 0.685 & 0.726 & 0.579 \\
control/IQP-no-entangle & 0.682 & 0.722 & 0.576 \\
quantum/IQP-kernel-SVM & 0.672 & 0.712 & 0.566 \\
quantum/Bures-RBF-SVM & 0.600 & 0.630 & 0.552 \\
quantum/HS-RBF-SVM & 0.595 & 0.621 & 0.549 \\
quantum/HS-overlap-SVM & 0.586 & 0.629 & 0.532 \\
quantum/Fidelity-SVM & 0.580 & 0.616 & 0.523 \\
\hline
\end{tabular}
\end{table}


%% ---------------------------------------------------------------- Table 6
\begin{table}[htbp]
\caption{\label{tab:cross}Pre-specified comparisons on both datasets, combined by
Fisher's method. Positive $\Delta$ favours the first-named pipeline. At $n=9$ the
two-sided signed-rank test cannot return $p<0.0039$, so the IV-2a entries marked
$\dagger$ sit at that floor: the test is saturated rather than merely significant. Both classical-versus-quantum contrasts hold in every one of the 9 IV-2a subjects.}
\begin{tabular}{@{}llccc@{}}
\hline
Comparison & Dataset & $\Delta$ acc & $p$ & Better \\
\hline
Best classical vs.\ best quantum & PhysioNet & $+0.0323$ & 0.002 & 24/30 \\
 & IV-2a & $+0.0662$ & 0.004$^\dagger$ & 9/9 \\
 & \textit{Fisher combined} & & \textbf{\ensuremath{<}0.001} & \\
\noalign{\smallskip}
Entanglement ablation & PhysioNet & $+0.0106$ & 0.309 & 17/30 \\
 & IV-2a & $+0.0100$ & 0.203 & 7/9 \\
 & \textit{Fisher combined} & & 0.237 & \\
\noalign{\smallskip}
Dimension-matched control & PhysioNet & $+0.0163$ & 0.173 & 18/30 \\
 & IV-2a & $+0.0178$ & 0.129 & 7/9 \\
 & \textit{Fisher combined} & & 0.107 & \\
\noalign{\smallskip}
\hline
\end{tabular}
\end{table}


%% -------------------------------------------------- Table: frame effect
\begin{table}[htbp]
\caption{\label{tab:frame}Effect of referring the density-matrix kernels to a
reference state. Each kernel is evaluated twice under an identical protocol,
differing only in whether the states are expressed in the sensor frame
(\ref{eq:density}) or relative to the training-set Fr\'echet mean
(\ref{eq:refstate}). $\Delta$ is the mean per-subject accuracy gain from the
reference frame, tested by paired Wilcoxon signed-rank across subjects;
$p_{\mathrm{Holm}}$ corrects for every row of this table as one family.
Every kernel improves on every dataset.}
\begin{tabular}{@{}llcccccc@{}}
\hline
Dataset & Kernel & Sensor & Reference & $\Delta$ & $p$ & $p_{\mathrm{Holm}}$ & Better \\
\hline
PhysioNet, 3\,q ($n=30$) & Fidelity & 0.523 & 0.630 & $\bf +0.1064$ & \ensuremath{<}0.001 & 0.005 & 23/30 \\
 & HS overlap & 0.532 & 0.621 & $\bf +0.0889$ & 0.002 & 0.024 & 21/30 \\
 & HS-RBF & 0.549 & 0.606 & $\bf +0.0570$ & 0.049 & 0.053 & 19/30 \\
 & Bures-RBF & 0.552 & 0.611 & $\bf +0.0585$ & 0.009 & 0.039 & 20/30 \\
 & QRE-RBF & 0.559 & 0.609 & $\bf +0.0504$ & 0.027 & 0.053 & 20/30 \\
\hline
PhysioNet, 4\,q ($n=30$) & Fidelity & 0.516 & 0.614 & $\bf +0.0983$ & 0.001 & 0.013 & 21/30 \\
 & HS overlap & 0.516 & 0.622 & $\bf +0.1062$ & \ensuremath{<}0.001 & 0.008 & 22/30 \\
 & HS-RBF & 0.543 & 0.619 & $\bf +0.0753$ & 0.004 & 0.039 & 20/30 \\
 & Bures-RBF & 0.543 & 0.626 & $\bf +0.0827$ & \ensuremath{<}0.001 & 0.008 & 23/30 \\
 & QRE-RBF & 0.541 & 0.624 & $\bf +0.0827$ & \ensuremath{<}0.001 & 0.004 & 26/30 \\
\hline
Cho2017, 3\,q ($n=52$) & Fidelity & 0.615 & 0.648 & $\bf +0.0332$ & \ensuremath{<}0.001 & \ensuremath{<}0.001 & 40/52 \\
 & HS overlap & 0.579 & 0.642 & $\bf +0.0633$ & \ensuremath{<}0.001 & \ensuremath{<}0.001 & 47/52 \\
 & HS-RBF & 0.600 & 0.637 & $\bf +0.0367$ & \ensuremath{<}0.001 & \ensuremath{<}0.001 & 42/52 \\
 & Bures-RBF & 0.627 & 0.639 & $\bf +0.0117$ & 0.009 & 0.039 & 33/52 \\
 & QRE-RBF & 0.629 & 0.641 & $\bf +0.0125$ & 0.002 & 0.023 & 36/52 \\
\hline
IV-2a, 3\,q ($n=9$) & Fidelity & 0.580 & 0.767 & $\bf +0.1873$ & 0.004 & 0.039 & \textbf{9/9} \\
 & HS overlap & 0.586 & 0.762 & $\bf +0.1760$ & 0.004 & 0.039 & \textbf{9/9} \\
 & HS-RBF & 0.595 & 0.763 & $\bf +0.1678$ & 0.004 & 0.039 & \textbf{9/9} \\
 & Bures-RBF & 0.600 & 0.765 & $\bf +0.1646$ & 0.004 & 0.039 & \textbf{9/9} \\
 & QRE-RBF & 0.599 & 0.764 & $\bf +0.1644$ & 0.004 & 0.039 & \textbf{9/9} \\
\hline
\end{tabular}
\end{table}


%% ------------------------------------------------------ Table: twin control
\begin{table}[htbp]
\caption{\label{tab:twin}The decisive control. Each quantum kernel is compared
against an SPD-manifold kernel used in the same support vector machine, on the
same covariances, with the same tuning budget, in the same reference frame:
only the metric differs. In every setting the twin is the affine-invariant
Riemannian kernel, the one classical kernel here that shares the
reference-frame quantum kernels' invariance under congruence. The
log-Euclidean kernel lacks that invariance and is shown for comparison only.
Columns give the best quantum kernel, the twin, the log-Euclidean kernel, the
range of the five quantum-minus-twin differences, and the smallest $p$
obtained by \emph{any} of the five quantum kernels against the twin.
Table~\ref{tab:seeds} repeats the PhysioNet three-qubit row under two further
partitions.}
\begin{tabular}{@{}lcccccc@{}}
\hline
Setting & $n$ & Quantum & Twin & Log-Eucl. & $\Delta$ (range over 5) & $\min p$ \\
\hline
PhysioNet, 3\,q & 30 & 0.630 & 0.621 & 0.612 & $[-0.0151, +0.0081]$ & 0.183 \\
PhysioNet, 4\,q & 30 & 0.626 & 0.625 & 0.623 & $[-0.0114, +0.0005]$ & 0.253 \\
PhysioNet, 5\,q, filter bank & 30 & 0.648 & 0.630 & 0.637 & $[+0.0069, +0.0180]$ & 0.066 \\
PhysioNet, transfer (LOSO) & 30 & 0.647 & 0.635 & 0.633 & $[+0.0000, +0.0119]$ & 0.072 \\
BCI IV-2a, 3\,q & 9 & 0.767 & 0.766 & 0.763 & $[-0.0033, +0.0016]$ & 0.426 \\
BCI IV-2a, cross-session & 9 & 0.760 & 0.755 & 0.755 & $[-0.0085, +0.0054]$ & 0.500 \\
Cho2017, 3\,q & 52 & 0.648 & 0.638 & 0.636 & $[-0.0008, +0.0100]$ & 0.001 \\
\hline
\end{tabular}
\end{table}


%% --------------------------------------------------- Table: equivalence
\begin{table}[htbp]
\caption{\label{tab:equiv}Equivalence of each quantum kernel to its classical
twin, by two one-sided tests on the paired per-subject differences. A
non-significant difference would only be an absence of evidence; TOST instead
takes a \emph{difference} as its null, so rejecting it supports equivalence.
The margin $m=0.02$ accuracy is not pre-registered. It
was chosen after the first within-subject analyses of PhysioNet and IV-2a, the
filter-bank analysis and the cross-subject analysis had been run, from the size
of the reference-frame effects they showed, and before every other analysis
reported here. It is left unchanged although some later effects are smaller
(the correction now spans \FrameEffectMin{} to \FrameEffectMax). CI is the
90\,\% interval, which
corresponds to TOST at $\alpha=0.05$; ``bound'' is the smallest margin at which
equivalence would hold, so a reader preferring a stricter margin can read the
answer off directly.}
\begin{tabular}{@{}llcccccc@{}}
\hline
Setting & Kernel & $n$ & $\Delta$ & 90\,\% CI & $p_{\mathrm{TOST}}$ &
Equiv. & Bound \\
\hline
PhysioNet, 3 qubits & Fidelity & 30 & $+0.0081$ & $[-0.0042, +0.0205]$ & 0.0568 & $\times$ & 0.0205 \\
 & HS overlap & 30 & $-0.0005$ & $[-0.0130, +0.0120]$ & 0.0066 & \checkmark & 0.0130 \\
 & HS-RBF & 30 & $-0.0151$ & $[-0.0317, +0.0016]$ & 0.3090 & $\times$ & 0.0317 \\
 & Bures-RBF & 30 & $-0.0109$ & $[-0.0255, +0.0038]$ & 0.1487 & $\times$ & 0.0255 \\
 & QRE-RBF & 30 & $-0.0121$ & $[-0.0267, +0.0025]$ & 0.1833 & $\times$ & 0.0267 \\
PhysioNet, 4 qubits & Fidelity & 30 & $-0.0114$ & $[-0.0263, +0.0036]$ & 0.1674 & $\times$ & 0.0263 \\
 & HS overlap & 30 & $-0.0037$ & $[-0.0163, +0.0089]$ & 0.0178 & \checkmark & 0.0163 \\
 & HS-RBF & 30 & $-0.0069$ & $[-0.0199, +0.0061]$ & 0.0493 & \checkmark & 0.0199 \\
 & Bures-RBF & 30 & $+0.0005$ & $[-0.0115, +0.0125]$ & 0.0050 & \checkmark & 0.0125 \\
 & QRE-RBF & 30 & $-0.0015$ & $[-0.0138, +0.0108]$ & 0.0081 & \checkmark & 0.0138 \\
PhysioNet, 5 qubits & Fidelity & 30 & $+0.0180$ & $[+0.0048, +0.0313]$ & 0.4010 & $\times$ & 0.0313 \\
 & HS overlap & 30 & $+0.0091$ & $[-0.0060, +0.0243]$ & 0.1167 & $\times$ & 0.0243 \\
 & HS-RBF & 30 & $+0.0069$ & $[-0.0075, +0.0214]$ & 0.0674 & $\times$ & 0.0214 \\
 & Bures-RBF & 30 & $+0.0086$ & $[-0.0027, +0.0200]$ & 0.0504 & $\times$ & 0.0200 \\
 & QRE-RBF & 30 & $+0.0084$ & $[-0.0042, +0.0210]$ & 0.0648 & $\times$ & 0.0210 \\
Transfer (LOSO) & Fidelity & 30 & $+0.0007$ & $[-0.0097, +0.0112]$ & 0.0020 & \checkmark & 0.0112 \\
 & HS overlap & 30 & $+0.0119$ & $[+0.0002, +0.0235]$ & 0.1231 & $\times$ & 0.0235 \\
 & HS-RBF & 30 & $+0.0000$ & $[-0.0154, +0.0154]$ & 0.0175 & \checkmark & 0.0154 \\
 & Bures-RBF & 30 & $+0.0052$ & $[-0.0080, +0.0183]$ & 0.0327 & \checkmark & 0.0183 \\
 & QRE-RBF & 30 & $+0.0067$ & $[-0.0068, +0.0201]$ & 0.0514 & $\times$ & 0.0201 \\
BCI IV-2a, 3 qubits & Fidelity & 9 & $+0.0016$ & $[-0.0049, +0.0080]$ & 0.0004 & \checkmark & 0.0080 \\
 & HS overlap & 9 & $-0.0033$ & $[-0.0109, +0.0043]$ & 0.0018 & \checkmark & 0.0109 \\
 & HS-RBF & 9 & $-0.0028$ & $[-0.0084, +0.0028]$ & 0.0002 & \checkmark & 0.0084 \\
 & Bures-RBF & 9 & $-0.0008$ & $[-0.0073, +0.0058]$ & 0.0003 & \checkmark & 0.0073 \\
 & QRE-RBF & 9 & $-0.0019$ & $[-0.0081, +0.0043]$ & 0.0003 & \checkmark & 0.0081 \\
IV-2a cross-session & Fidelity & 9 & $+0.0042$ & $[-0.0057, +0.0142]$ & 0.0094 & \checkmark & 0.0142 \\
 & HS overlap & 9 & $-0.0004$ & $[-0.0124, +0.0117]$ & 0.0082 & \checkmark & 0.0124 \\
 & HS-RBF & 9 & $-0.0085$ & $[-0.0281, +0.0111]$ & 0.1531 & $\times$ & 0.0281 \\
 & Bures-RBF & 9 & $+0.0054$ & $[-0.0138, +0.0246]$ & 0.0975 & $\times$ & 0.0246 \\
 & QRE-RBF & 9 & $+0.0054$ & $[-0.0144, +0.0252]$ & 0.1034 & $\times$ & 0.0252 \\
Cho2017, 3 qubits & Fidelity & 52 & $+0.0100$ & $[+0.0047, +0.0153]$ & 0.0013 & \checkmark & 0.0153 \\
 & HS overlap & 52 & $+0.0041$ & $[-0.0010, +0.0093]$ & 0.0000 & \checkmark & 0.0093 \\
 & HS-RBF & 52 & $-0.0008$ & $[-0.0082, +0.0066]$ & 0.0000 & \checkmark & 0.0082 \\
 & Bures-RBF & 52 & $+0.0010$ & $[-0.0063, +0.0083]$ & 0.0000 & \checkmark & 0.0083 \\
 & QRE-RBF & 52 & $+0.0034$ & $[-0.0036, +0.0104]$ & 0.0001 & \checkmark & 0.0104 \\
\hline
\end{tabular}
\end{table}


%% --------------------------------------------------- Table: transfer
\begin{table}[htbp]
\caption{\label{tab:transfer}Cross-subject transfer, leave-one-subject-out:
each model is trained on the pooled trials of every other subject.
Hyperparameters are selected by subject-grouped inner cross-validation on the
training subjects only. In the reference frame each subject is whitened by its
own Fr\'echet mean, which uses no labels and is therefore legitimate
unsupervised adaptation for the held-out subject. $\Delta$ is the gain from the
reference frame, paired by held-out subject. Every method improves; in the
reference frame no method is distinguishable from any other.}
\begin{tabular}{@{}lccccc@{}}
\hline
Pipeline & Sensor & Reference & $\Delta$ & $p$ & Better \\
\hline
quantum/HS-overlap & 0.552 & 0.647 & $+0.0948$ & \ensuremath{<}0.001 & 24/30 \\
classical/TS+LR & 0.604 & 0.644 & $+0.0400$ & 0.010 & 18/30 \\
quantum/QRE-RBF & 0.576 & 0.641 & $+0.0659$ & 0.002 & 21/30 \\
quantum/Bures-RBF & 0.561 & 0.640 & $+0.0785$ & \ensuremath{<}0.001 & 22/30 \\
quantum/Fidelity & 0.559 & 0.636 & $+0.0763$ & \ensuremath{<}0.001 & 21/30 \\
control/riemann-kernel-SVM & 0.587 & 0.635 & $+0.0474$ & 0.004 & 19/30 \\
quantum/HS-RBF & 0.564 & 0.635 & $+0.0711$ & 0.002 & 20/30 \\
control/logeuclid-kernel-SVM & 0.576 & 0.633 & $+0.0570$ & 0.001 & 20/30 \\
classical/MDM & 0.553 & 0.633 & $+0.0800$ & \ensuremath{<}0.001 & 19/30 \\
\hline
\end{tabular}
\end{table}


%% ----------------------------------------------- Table: cross-session
\begin{table}[htbp]
\caption{\label{tab:crosssession}Cross-session transfer within subject on
IV-2a. For each subject and in both directions, each model is trained on one
recording session and tested on the other; the two directions are averaged
per subject before any statistic is computed, so $n$ is the number of
subjects. Hyperparameters are chosen by stratified inner cross-validation
inside the training session alone. In the reference frame each session is
whitened by its own label-free Fr\'echet mean. $\Delta$ is the gain from the
reference frame, paired by subject and tested by Wilcoxon signed-rank.}
\begin{tabular}{@{}lccccc@{}}
\hline
Pipeline & Sensor & Reference & $\Delta$ & $p$ & Better \\
\hline
quantum/QRE-RBF & 0.576 & 0.760 & $+0.1844$ & 0.004 & \textbf{9/9} \\
quantum/Bures-RBF & 0.573 & 0.760 & $+0.1867$ & 0.004 & \textbf{9/9} \\
quantum/Fidelity & 0.542 & 0.759 & $+0.2172$ & 0.004 & \textbf{9/9} \\
control/riemann-kernel-SVM & 0.734 & 0.755 & $+0.0204$ & 0.012 & 8/9 \\
control/logeuclid-kernel-SVM & 0.724 & 0.755 & $+0.0309$ & 0.020 & 8/9 \\
quantum/HS-overlap & 0.560 & 0.754 & $+0.1944$ & 0.008 & 8/9 \\
quantum/HS-RBF & 0.578 & 0.746 & $+0.1686$ & 0.004 & \textbf{9/9} \\
classical/TS+LR & 0.717 & 0.738 & $+0.0216$ & 0.039 & 8/9 \\
classical/MDM & 0.701 & 0.732 & $+0.0312$ & 0.023 & 7/9 \\
\hline
\end{tabular}
\end{table}


%% ------------------------------------------------------ Table: seeds
\begin{table}[htbp]
\caption{\label{tab:seeds}The PhysioNet three-qubit comparison under three
different outer cross-validation partitions (seeds). Each row repeats the
full extended suite with new fold assignments; the inner search, the reference
state and every other element are unchanged. Columns give the frame effect
(range over the five kernels), the best reference-frame quantum kernel and
the classical twin, the quantum-minus-twin difference (range over the five
kernels), the smallest $p$ against the twin, and how many of the five
comparisons pass two one-sided tests at the pre-specified $\pm0.02$
margin, with the largest bound. The frame effect and the uncontrolled
reversal reproduce in every partition; the twin comparison moves against the
bandwidth-parameterised quantum kernels under the other two partitions and
never in their favour (see text).}
\begin{tabular}{@{}lcccccc@{}}
\hline
Seed & Frame $\Delta$ & Quantum & Twin & Quantum $-$ twin & $\min p$ &
TOST \\
\hline
0 & $[+0.050, +0.106]$ & 0.630 & 0.621 & $[-0.0151, +0.0081]$ & 0.183 & 1/5 ($\le0.032$) \\
1 & $[+0.039, +0.104]$ & 0.613 & 0.609 & $[-0.0220, +0.0040]$ & 0.030 & 2/5 ($\le0.039$) \\
2 & $[+0.032, +0.107]$ & 0.618 & 0.616 & $[-0.0284, +0.0022]$ & \ensuremath{<}0.001 & 1/5 ($\le0.042$) \\
\hline
\end{tabular}
\end{table}


\subsection{Quantum kernels do not outperform tuned classical baselines}

Table~\ref{tab:main} reports within-subject performance for all \NPipelines{}
pipelines; Figure~\ref{fig:benchmark} shows the same data with per-subject spread.
Quantum kernels occupy five of the bottom six positions. Absolute accuracies
are modest because EEGMMIDB contains many near-chance subjects, which is
precisely why paired testing is required. They are also in line with the
field: the MOABB reproducibility benchmark \cite{chevallier2024} reports
65.8--68.2\,\% for CSP and tangent-space pipelines on this dataset and
paradigm using the full 64-channel montage, against which our
\BestClassicalAcc{} on eight channels is the expected reduction, not a sign of
a mistuned baseline.

The primary pre-specified comparison, best classical versus best quantum, is
significant: \BestClassicalName{} exceeded \BestQuantumName{} by
$\Delta=\PrimaryDelta$ ($p\PrimaryP$, $d_z=\PrimaryDz$), and was superior in
\PrimaryBetter{} subjects (Table~\ref{tab:key}).

Against the Riemannian tangent-space reference (Table~\ref{tab:tests},
Figure~\ref{fig:paired}), \NSurviveHolm{} of \NFamilyTests{} comparisons survived
Holm correction. Only two reached $p<0.05$ uncorrected, and in both the quantum
model was \emph{worse}. No quantum kernel exceeded the reference at any
significance level.

\begin{figure}[htbp]
\centering
\includegraphics[width=\textwidth]{fig2_benchmark.pdf}
\caption{\label{fig:benchmark}Within-subject accuracy by pipeline
($n=\NSubjects$ subjects, nested cross-validation). Bars are means over
subjects; grey dots are individual subject means. The dashed line marks chance.
Every pipeline, classical baselines included, received the same inner
hyperparameter search.}
\end{figure}

\begin{figure}[htbp]
\centering
\includegraphics[width=\textwidth]{fig3_paired_differences.pdf}
\caption{\label{fig:paired}Paired per-subject accuracy differences against the
Riemannian tangent-space reference. Dots are individual subjects; diamonds are
means. Points left of zero indicate worse performance than the reference.}
\end{figure}

\subsection{The controls: entanglement and dimension matching}

Both control ablations point in the direction predicted by \cite{bowles2024},
but neither reaches significance on accuracy (Table~\ref{tab:key}). Removing
entanglement changed accuracy by $\AblationDelta$ ($p\AblationP$, superior in
\AblationBetter{} subjects); the dimension-matched linear SVM exceeded the IQP
circuit kernel by $\DimDelta$ ($p\DimP$).

We regard the honest reading as follows: deleting the quantum resource did not
harm performance and numerically helped, but accuracy alone, measured on a
noisy and low-SNR signal, does not have the power to establish the effect. The
next section shows that the same effect is unambiguous when measured on the
kernel itself.

\subsection{Kernel concentration, and why entanglement makes it worse}
\label{sec:concentration}

Raw overlap kernels are severely concentrated on EEG. Off-diagonal Gram entries
had mean $0.991\pm0.010$ (Hilbert--Schmidt) and $0.988\pm0.010$ (fidelity):
every pair of trials appears $\approx99\%$ identical, the Gram matrix is
effectively rank-one, and the classifier has nothing to separate. This is the
pathology proved by Thanasilp \textit{et al} \cite{thanasilp2024}, here reproduced on
real biological data rather than synthetic constructions.

Exploiting (\ref{eq:density}), we swept register size by varying the channel
count over powers of two ($4/8/16/32/64$ channels $=2$--$6$ qubits) with no
change of method (Table~\ref{tab:concentration}, Figure~\ref{fig:concentration}). Two
findings emerge.

First, \textbf{entanglement accelerates the collapse}. The entangled IQP kernel
loses variance by a factor of $0.428$ per qubit, more than halving with each
added qubit, while the identical circuit with entanglers deleted decays at
$0.716$. In log-slope the entangled kernel concentrates $2.5\times$ faster
($-0.850$ versus $-0.334$ per qubit), a $29\times$ versus $3.7\times$ loss
across the range. This supplies a mechanism for the observation of
\cite{bowles2024} that removing entanglement often improves QML models:
entanglement drives embedded states toward mutual orthogonality, pushing the
Gram matrix toward the identity, a memorising, non-generalising model.

Second, \textbf{the two kernel families concentrate for different reasons}. The
raw density-matrix kernels do not follow the qubit-count law at all; their
variance \emph{increases} with dimension (factors $1.45$--$1.52$). They are
already saturated at two qubits because EEG covariance matrices are
intrinsically similar to one another, a property of the data, not of the
register. The practical corollary inverts standard practice: the density-matrix
route should use \emph{more} channels, not fewer.

\begin{figure}[htbp]
\centering
\includegraphics[width=\textwidth]{fig1_concentration.pdf}
\caption{\label{fig:concentration}Kernel variance versus register size on real
EEG ($n=10$ subjects). Lower is worse: a kernel whose variance collapses cannot
separate trials. Line style encodes the kernel family. The entangled circuit
kernel (orange, dashed) decays fastest.}
\end{figure}

\subsection{Replication on IV-2a: more data widens the gap}
\label{sec:crossdataset}

Repeating the entire benchmark on BCI Competition IV-2a reproduces the
ordering and sharpens it (Table~\ref{tab:bci}, Figure~\ref{fig:cross}). The best classical
pipeline reached \BciBestClassicalAcc{} against \BciBestQuantumAcc{} for the
best quantum kernel and \BciWorstQuantumAcc{} for the worst, a
classical-versus-quantum gap of \GapBci{} where PhysioNet gave \GapPhysio{},
and a gap to the weakest quantum kernel of \GapWorstBci{} against
\GapWorstPhysio{}.

The direction of that change is the point. Classical methods gained roughly
0.15 accuracy from the six-fold increase in trials; the density-matrix kernels
gained about 0.05. \textbf{More and cleaner data made the quantum kernels
relatively worse, not better}, which is the opposite of what a data-starvation
explanation predicts, and exactly what section~\ref{sec:concentration} implies: a
kernel whose off-diagonal entries all sit near unity gains no additional
structure from additional trials, while the classical methods gain a great
deal.

The ranking is also stable. Spearman rank correlation between the two
orderings is \SpearmanRho{} across all \NPipelines{} pipelines, and the four
lowest positions are occupied by the same four quantum kernels on both
datasets with zero rank shift. This is a property of the methods, not of the
dataset.

Statistically (Table~\ref{tab:cross}), on IV-2a the best classical pipeline beat the
best quantum kernel by $\BciPrimaryDelta$ ($p\BciPrimaryP$,
$d_z=\BciPrimaryDz$) in \BciPrimaryBetter{} subjects, and beat the
density-matrix kernel by $\BciDensDelta$ ($d_z=\BciDensDz$) in
\BciDensBetter{} subjects. Combining the two independent datasets by Fisher's
method gives $p\FisherPrimaryP$ and $p\FisherDensP$ respectively.

One caveat must be stated precisely rather than left to the reader. With $n$
pairs the two-sided signed-rank test cannot return a $p$ below $2^{1-n}$; at
$n=\BciNSubjects$ that floor is \BciWilcoxonFloor{}, so a Holm correction over
\NFamilyTests{} comparisons cannot fall below \BciHolmFloor{} whatever the
effect size. Ten of the fourteen IV-2a family comparisons sit exactly at that
floor. Reporting them as "not significant after correction" would misdescribe
a saturated test as a null result, which is why the pre-specified comparisons
in Table~\ref{tab:cross}, which require no family correction, carry the inference.

\begin{figure}[htbp]
\centering
\includegraphics[width=\textwidth]{fig4_crossdataset.pdf}
\caption{\label{fig:cross}The same \NPipelines{} pipelines on two datasets at
opposite corners of the design space. Identical channels, protocol, tuning
budget and code; only the data differs. Line slope is the quantity of
interest: classical and control pipelines rise steeply with the extra trials
while the four density-matrix and overlap kernels stay low. The dashed line
marks chance.}
\end{figure}

\subsection{The sensor frame is a confound, and correcting it reverses the
result}
\label{sec:refresults}

Every result above evaluates the density-matrix kernels in the sensor frame.
Section~\ref{sec:reference} argued that this is not neutral: the quantum kernels are
invariant only under orthogonal congruence, while the tangent-space, MDM and
CSP baselines they are being compared against are invariant under the full
congruence group that EEG's nuisances generate. We now re-run the entire
benchmark with the density-matrix kernels expressed relative to a training-set
reference state (\ref{eq:refstate}), changing nothing else. The core
pipelines reproduce their published values to $\le 0.0001$, so the two suites
are directly comparable.

Three consequences follow, and the third is the one that matters.

\emph{First, concentration is largely an artefact of the frame.} Measured over
\GramNSubjects{} subjects, the off-diagonal variance of the Gram matrix rises
by a factor of \GramGainMin--\GramGainMax{} on moving to the reference frame.
The sensor-frame Hilbert--Schmidt kernel has off-diagonal mean
\GramSensorMean{} and standard deviation \GramSensorStd; a nuisance congruence
of the kind an electrode-gain change produces shifts individual kernel entries
by an order of magnitude more than that entire discriminative spread. The
kernel was measuring the frame, not the trial.

The register sweep of section~\ref{sec:concentration} was made in the sensor
frame, and its second finding, that the density-matrix kernels are saturated
by the data rather than by the register, has to be re-examined now that most
of that saturation is known to be an artefact of the frame. Repeating the
sweep in both frames (table~\ref{tab:sweep}, figure~\ref{fig:crosssession}(b))
gives a cleaner answer. In the reference frame no kernel's Gram variance falls
below its three-qubit value at any larger register, and at six qubits it
stands at \SweepRefRatioMin--\SweepRefRatioMax{} times that value; the
exponential collapse with qubit count of \cite{thanasilp2024} is absent from
this family in either frame. The gain from recentring narrows from
\SweepGainLoMin--\SweepGainLoMax$\times$ at three qubits to
\SweepGainHiMin--\SweepGainHiMax$\times$ at six, but that is because the
sensor-frame variance itself rises
\SweepSensorRatioMin--\SweepSensorRatioMax$\times$ over the same range, not
because the reference frame deteriorates. The corollary of
section~\ref{sec:concentration}, that the density-matrix route tolerates wider
registers, therefore survives the correction, with one qualification: the
reference-frame variance is not strictly monotone, three of the four kernels
dipping by up to \SweepRefDipMax\,\% between five and six qubits.

\emph{Second, accuracy improves substantially and consistently}
(table~\ref{tab:frame}). On PhysioNet every kernel improves, by \PhysFrameMin{} to
\PhysFrameMax{} ($n=\PhysFrameN$, all $p \le \PhysFrameMaxP$). On IV-2a the
same corrections are worth \BciFrameMin{} to \BciFrameMax{}, and every kernel
improves in \emph{every} subject. That the effect is two to three times larger
on IV-2a is expected rather than surprising: the reference state is a
Fr\'echet mean estimated from training covariances, and \BciNTrials{} trials
per subject estimate it far better than \NTrials{} do.

This also revises the interpretation of section~\ref{sec:crossdataset}. There, more
data widened the classical--quantum gap, which we read as evidence that the
quantum kernels could not exploit additional trials. The narrower and correct
reading is that a kernel with the wrong invariance cannot exploit them. Given
the right frame, the extra data is not merely usable but more valuable to the
quantum kernels than to the classical baselines.

The headline comparison consequently reverses. Against \HeadBestClassical, the
best sensor-frame quantum kernel loses by \HeadSensorDelta{}
($p\HeadSensorP$, classical better in \HeadSensorBetter{} subjects); the best
reference-frame kernel, \HeadRefKernel, instead leads by
\HeadRefLead{} ($p\HeadRefP$, classical better in only
\HeadRefBetter). This does not reach conventional significance, and we do not
claim it as an effect. We dwell on it because it is the size and shape of
result that this literature routinely reports as a positive finding, and
because everything needed to present it as one --- two datasets, a consistent
direction, a larger margin on the better-powered data --- is in hand at this
point in the analysis.

\emph{Third, it is not a quantum advantage, and the control that establishes
this is the central result of this paper.} A reference-frame quantum kernel
differs from a classical Riemannian kernel in exactly one respect: the metric.
Both consume the same covariances, both are used as precomputed kernels in the
same support vector machine, both receive the same tuning budget, and after
(\ref{eq:refstate}) both are affine-invariant. Table~\ref{tab:twin} reports that
comparison in \EquivNSettings{} settings: \EquivNDatasets{} datasets,
\EquivNRegisters{} register sizes, and within-subject, cross-subject and
cross-session evaluation.

In none of them, at the primary partition, is any quantum kernel
distinguishable from its classical twin (figure~\ref{fig:twin}).
On PhysioNet the differences span \TwinDeltaMin{} to \TwinDeltaMax{} with a
smallest $p\TwinMinP$ in that setting; the twin itself, \TwinName,
reaches \TwinAcc{} and is the second-best pipeline of the whole extended suite.
The apparent quantum advantage is therefore attributable to two things, neither
of them quantum: expressing the states in the reference frame, and using an
SPD-manifold kernel in a support vector machine rather than a tangent-space
projection.

A non-significant difference is only an absence of evidence, and the claim we
are making deserves better than that, so we test equivalence directly
(table~\ref{tab:equiv}). Two one-sided tests on the paired differences, with a
margin of $\pm\EquivMargin$ accuracy fixed in advance, establish equivalence in
\EquivPass{} of \EquivN{} comparisons; across all \EquivN{} the agreement is
within $\pm\EquivWorstBound$. The evidence is strongest exactly where it should
be, on the better-powered dataset: on IV-2a all \EquivBciN{} comparisons are
equivalent, with the widest interval reaching only $\EquivBciWorst$.

Two honest qualifications. First, the \EquivFailQuantum{} comparisons that miss
the $\pm\EquivMargin$ margin in the quantum kernel's favour are balanced by
\EquivFailClassical{} that miss it in the twin's, which is the pattern expected
of noise about zero rather than of a suppressed effect. Second, exactly
\EquivNZeroExcluded{} of the \EquivN{} intervals excludes zero at all:
\EquivZeroExclName, at $\EquivZeroExclDelta$. We do not read that as an
advantage --- it is one comparison of \EquivN{}, it does not survive the
family, and its interval still lies inside $\pm\EquivWorstBound$ --- but it is
the single place in this study where a quantum kernel is nominally ahead of a
metric-matched classical control, and a reader is entitled to see it named.

\emph{Robustness to the partition.} Every number above comes from one outer
cross-validation partition. Table~\ref{tab:seeds} repeats the PhysioNet
three-qubit suite under two more. The frame effect reproduces
(\SeedFrameMin{} to \SeedFrameMax{} over all \SeedN{} partitions), and so
does the uncontrolled reversal: the best reference-frame kernel leads the
best classical baseline by \SeedHeadRefMin{} to \SeedHeadRefMax{} in
\SeedHeadReversals{} of \SeedN{} partitions. The twin comparison is less
forgiving than the primary partition suggests, and in one direction only.
The parameter-free fidelity and Hilbert--Schmidt kernels stay tied to the
twin under every partition (\SeedOverlapDeltaMin{} to
\SeedOverlapDeltaMax, smallest $p\SeedOverlapMinP$), but the three
bandwidth-parameterised kernels, which carry an extra hyperparameter to be
tuned on \NTrials{} trials, fall behind it by up to \SeedRbfWorst{} with a
smallest $p\SeedRbfMinP$; \SeedSigCount{} of the \SeedSigN{} comparisons
across the three partitions are significant at $0.05$, \SeedSigFavourTwin{}
in the twin's favour and \SeedSigFavourQuantum{} in the quantum kernel's.
The equivalence bound for this dataset widens accordingly, from
\EquivWorstBound{} at the primary partition to \SeedWorstBound{} over all
three. We therefore state the PhysioNet result as an inequality rather than
an equivalence: the quantum kernels match the classical twin at best, and a
bandwidth to tune on \NTrials{} trials costs them a few points against it.
Nothing in the repetitions moves in the quantum direction.

We stress the counterfactual. Without this control, the same experiments
support the claim that quantum kernels outperform tuned classical baselines at
$p\HeadRefP$, on two datasets, with the gap widening on the better one. That
claim would have been false. It is also, we suggest, the mechanism by which
favourable results are obtained elsewhere in this literature: a quantum kernel
that operates on whitened or aligned covariances, compared against a classical
baseline that does not, will show an advantage of roughly the size reported in
table~\ref{tab:frame} that has nothing to do with quantum structure.

\begin{figure}[htbp]
\centering
\includegraphics[width=\textwidth]{fig7_frame.pdf}
\caption{\label{fig:frame}Accuracy in the sensor frame versus the reference
frame. Thin lines are individual subjects, shown for the Fidelity kernel only.
Every kernel improves on both datasets; on IV-2a every kernel improves in every
subject. The correction is worth two to three times more on IV-2a, where the
reference state is estimated from \BciNTrials{} rather than \NTrials{} trials
per subject.}
\end{figure}

\begin{figure}[htbp]
\centering
\includegraphics[width=\textwidth]{fig8_twin.pdf}
\caption{\label{fig:twin}The decisive control. Each quantum kernel is compared
against a classical SPD-manifold kernel \cite{barachant2013} used in the
same support vector
machine, on the same covariances, with the same tuning budget, in the same
reference frame; only the metric differs. Points are mean paired differences
with 95\% confidence intervals for the difference; table~\ref{tab:equiv}
gives the 90\% intervals that correspond to the equivalence test, which are
narrower. Every interval here crosses zero in all \EquivNSettings{} settings
at the primary partition (table~\ref{tab:seeds} for two more). Note the axis
range: the entire plot spans eight accuracy points, against a frame effect
(tables~\ref{tab:frame} and \ref{tab:crosssession}) of up to $\CsFrameQMax$.}
\end{figure}

\subsection{Filter-bank baselines and larger registers}
\label{sec:filterbank}

Since the reference-frame kernels reach the top of the table, the relevant
question becomes whether they do so against the strongest classical family
rather than against CSP and tangent-space alone. We therefore repeated the
comparison with a filter bank of four sub-bands spanning the same passband,
in the manner of filter-bank CSP \cite{ang2008},
giving filter-bank CSP and filter-bank tangent-space baselines. The same
construction enlarges the quantum register at no methodological cost: four
sub-bands over eight channels is a $32$-dimensional covariance, hence a single
density matrix on five qubits rather than three.

Parity survives. The best quantum kernel reaches \FbcspBestAcc{} against
\FbcspAcc{} for filter-bank CSP, a difference of \FbcspDelta{}
($p\FbcspP$), and the classical-twin control replicates at this register size
as well (table~\ref{tab:twin}). The frame effect is larger here than at three
qubits, \FbFrameMin{} to \FbFrameMax; in the sensor frame the five-qubit
kernels fall to chance and are the only comparisons in that family to survive
Holm correction, in the unfavourable direction.

The other route to a larger register is more channels. Running the same
suite on sixteen sensorimotor channels, a four-qubit register, gives the same
picture (tables~\ref{tab:frame} and \ref{tab:twin}): the frame is worth
\FourQFrameMin{} to \FourQFrameMax{} ($n=\FourQFrameN$, all
$p\le\FourQFrameMaxP$), the best reference-frame kernel, \FourQBestKernel,
reaches \FourQBestAcc{} against \FourQTwinAcc{} for the twin, and the five
differences span \FourQTwinDeltaMin{} to \FourQTwinDeltaMax{} (smallest
$p\FourQTwinMinP$). This register is also where the uncontrolled comparison
is most seductive. Against \FourQBestClassical{} the best reference-frame
quantum kernel leads by \FourQHeadRefDelta{} ($p\FourQHeadRefP$, better in
\FourQHeadRefBetter{} subjects), which unlike the three-qubit case
(section~\ref{sec:refresults}) is conventionally significant; the classical
twin leads the same baseline by \FourQTwinHeadDelta{} ($p\FourQTwinHeadP$).
Without the twin, four qubits would have supplied a significant quantum
advantage on thirty subjects.

\subsection{Cross-subject transfer}
\label{sec:transfer}

Within-subject accuracy is not the field's bottleneck; transfer to a new user
is. It is also the setting in which a different geometry is most often
conjectured to help. Table~\ref{tab:transfer} and figure~\ref{fig:transfer}(a)
report leave-one-subject-out evaluation over \TransferN{} subjects, training on
the pooled trials of all others.

The invariance argument makes a sharp prediction here. Inter-subject shift is
dominated by congruence, so in the sensor frame the quantum distances must be
worse than the affine-invariant Riemannian one, and in the reference frame they
can at best equal it: parity is the ceiling the algebra permits, and no
experiment can exceed it. Both halves hold. Every method gains from
recentring, the non-invariant quantum kernels gaining most; and in the
reference frame the best method, \TransferBest{} at \TransferBestAcc{}, is not
distinguishable from any other. All nine geometries fall within
\TransferSpread{} of one another, with no comparison significant even before
correction. This is a genuine null rather than an underpowered one: at
$n=\TransferN$ the two-sided Wilcoxon signed-rank test could have returned $p$
as small as $\TransferFloor$.

\subsection{Cross-session transfer}
\label{sec:crosssession}

Cross-subject shift is dominated by congruence, so the parity just reported is
what the algebra predicts. The complementary case is transfer across sessions
within a subject: same head, same cap, similar impedances, so the congruence
component is milder and what remains, drift in conditioning, effective rank
and spectral shape, is exactly the residual that the invariance argument of
section~\ref{sec:reference} says nothing about. It is where a quantum
divergence, which weights the eigenvalue spectrum differently from the
affine-invariant metric, would separate from its classical twin if it were
going to anywhere, and we had named it in advance as the most promising place
to look. IV-2a provides it: two recording sessions of \CsTrialsPerSession{}
trials per subject. For each subject and in both directions we train on one
session and test on the other (table~\ref{tab:crosssession},
figure~\ref{fig:crosssession}(a)); hyperparameters are tuned inside the
training session only, and in the reference frame each session is whitened by
its own label-free mean, which for the test session is the unsupervised
adaptation a deployed decoder can legitimately perform.

Two things happen. First, the frame matters more here than anywhere else in
the study. In the sensor frame the quantum kernels reach only \CsSensorQMin{}
to \CsSensorQMax, close to chance, while the classical baselines, which do not
see the session change as a change of frame, reach \CsSensorClMin{} to
\CsSensorClMax. Recentring is worth \CsFrameQMin{} to \CsFrameQMax{} to the
quantum kernels (all $p\le\CsFrameQMaxP$) and \CsFrameClMin{} to
\CsFrameClMax{} to the classical ones. Session-to-session change evidently has
a large congruence component after all, whatever its residual, and a kernel
without the invariance is defenceless against it.

Second, once the frame is corrected the geometries tie again. The best
pipeline is \CsBestKernel{} at \CsBestAcc; the classical twin reaches
\CsTwinAcc{} and the best conventional baseline, \CsBestClassical,
\CsBestClassicalAcc. Against the twin the five quantum kernels differ by
\CsTwinDeltaMin{} to \CsTwinDeltaMax, and the smallest $p$ among the five
comparisons, $p\CsTwinMinP$, is the largest such value in any setting of the
study; the five quantum kernels and the two classical twins lie within
\CsSpread{} of one another. The equivalence bounds
(table~\ref{tab:equiv}) are wider than in the within-subject IV-2a comparison,
as they must be with one accuracy per subject rather than fifteen folds, but
the largest is \EquivCsWorst, inside the study-wide $\pm\EquivWorstBound$. The
setting we identified in advance as the quantum divergences' best chance is
therefore one more in which they match a classical kernel that differs only in
metric. Either the non-congruence part of session shift is too small to matter
at this register, or it is not something the Bures, fidelity or
relative-entropy weightings exploit better than the affine-invariant one.

\begin{figure}[htbp]
\centering
\includegraphics[width=\textwidth]{fig10_crosssession_sweep.pdf}
\caption{\label{fig:crosssession}(a) Cross-session transfer within subject on
IV-2a, both directions averaged. In the sensor frame the quantum kernels sit
near chance while the classical baselines do not; whitening each session by
its own label-free mean closes the gap and leaves every SPD-kernel geometry
tied. (b) Off-diagonal Gram variance versus register size in both frames on
PhysioNet ($n=\SweepN$). Solid, reference frame; dashed, sensor frame. In the
reference frame no kernel falls below its three-qubit value at any larger
register.}
\end{figure}

\subsection{Finite-shot estimation}
\label{sec:shots}

All results so far assume noiseless, infinite-shot simulation. Since
$\mathrm{tr}(\rho\sigma)$ is the SWAP-test observable, $S$ shots give an
unbiased estimate with variance $(1-k^2)/S$, which we simulate by sampling
each unordered Gram entry binomially and projecting back to the positive
semi-definite cone (table~\ref{tab:shots} and figure~\ref{fig:transfer}(b)).
One departure from the main protocol should be stated: here the reference
state is estimated from all of a subject's trials rather than the training
split alone. It is label-free, and the question is estimation cost rather than
generalisation, so shot levels are comparable with one another; the absolute
values are not comparable with the nested-CV benchmark and are not used that
way.

The variance gain reported above does \emph{not} translate into a
proportionate saving in shots; the reference frame requires more shots to
approach its own ceiling, because that ceiling is higher and the distinctions
to be resolved are correspondingly finer. Normalising each frame by its own
ceiling is in fact misleading, since the sensor frame converges quickly only
because it converges to near-chance. The meaningful comparison is absolute:
from $\ShotCrossover$ shots per Gram entry upwards, the reference frame exceeds
what the sensor frame attains with unlimited shots. Below $10^3$ shots both
collapse to chance, so the frame raises the ceiling a well-resourced estimator
can reach rather than rescuing a starved one.

Near-ceiling accuracy nonetheless requires $10^5$--$10^6$ shots per Gram entry.
A \NTrials-trial subject has of order $10^3$ unordered pairs, so a single Gram
matrix costs $10^8$--$10^9$ shots, and the IV-2a trial count raises that by a
further factor of about forty. The reference frame makes these kernels work;
it does not make them cheap.

\begin{figure}[htbp]
\centering
\includegraphics[width=\textwidth]{fig9_transfer_shots.pdf}
\caption{\label{fig:transfer}(a) Cross-subject transfer, leave-one-subject-out.
Recentring helps every method and helps the quantum kernels (Q) most, but in
the reference frame all nine geometries are indistinguishable. (b) Accuracy
against shot budget. The reference frame needs more shots to approach its own
ceiling because that ceiling is higher, yet from $\ShotCrossover$ shots per
Gram entry it already exceeds what the sensor frame attains with unlimited
shots.}
\end{figure}

\subsection{Computational cost}

The cost asymmetry survives the frame correction, which is the comparison
that matters now that accuracy is at parity. In the sensor frame the gap is
stark but partly beside the point, since the quantum kernels were also losing:
\FastestName{} reached its accuracy in \FastestSec~s per subject while
\SlowestName{} required \SlowestSec~s, a factor of \SpeedRatio. In the
reference frame the like-for-like comparison is against the classical twin,
which delivers the same accuracy. On IV-2a the twin costs \BciCostTwinSec~s
per subject; \BciCostSlowName{} costs \BciCostSlowSec~s, a factor of
\BciCostRatio, for an accuracy difference table~\ref{tab:equiv} bounds at
under $\pm\EquivBciWorst$. On PhysioNet the ratio is \PhysCostRatio. The
quantum kernel is not merely no better; it is the same answer at several times
the price.

All of this is measured under noiseless, infinite-shot simulation. On
hardware, estimating $\mathrm{tr}(\rho\sigma)$ to precision $\epsilon$ costs
$O(\epsilon^{-2})$ shots per Gram entry; section~\ref{sec:shots} puts the
resulting budget at $10^5$--$10^6$ shots per entry for near-ceiling accuracy,
which the classical twin does not pay at all.

%% =====================================================================
\section{Discussion}
%% =====================================================================

\subsection{Interpretation}

Under matched tuning, matched features and paired statistics, quantum kernels
did not improve motor-imagery decoding. That conclusion is consistent with the
broader QML benchmarking literature \cite{bowles2024}, but the route by which
we reach it differs from a bare negative result in a way we think matters, and
it passes through a stage at which the opposite conclusion looked correct.

Evaluated in the sensor frame, the density-matrix kernels, the most principled
construction here, requiring no dimensionality reduction and mapping onto a
hardware-native SWAP test, performed \emph{worst} of all pipelines, and
section~\ref{sec:concentration} appeared to explain why: EEG covariances are mutually
so similar that their quantum overlaps saturate near unity. That explanation is
incomplete. Concentration is a symptom of a deeper mismatch, and it is largely
removable. Uhlmann fidelity, the Bures distance, the Hilbert--Schmidt overlap
and the quantum relative entropy are invariant under \emph{unitary} conjugation;
EEG's nuisances act by \emph{congruence}. The affine-invariant Riemannian metric
is invariant under that whole group, which is the structural reason tangent-space
decoding is the classical state of the art. Comparing the two families in the
sensor frame therefore measures, in part, a difference of invariance groups
rather than a difference of geometries.

Correcting this (\ref{eq:refstate}) relieves concentration by
\GramGainMin--\GramGainMax$\times$, lifts every kernel on both datasets, and
reverses the sign of the headline comparison. Taken alone, that is a
publishable claim of quantum advantage at $p\HeadRefP$, replicated on a second
dataset, with the effect \emph{growing} on the better-powered one. We think it
is worth being explicit that this is what our own experiments supported at that
point, because the only thing standing between it and publication was a control
that is not currently standard.

That control, an SPD kernel in the same classifier, same data, same frame,
differing only in the metric, removes the claim entirely. Across \EquivNSettings{} settings
no quantum kernel exceeds its classical twin; the twin is itself among the
strongest pipelines we evaluated, and under other partitions of PhysioNet it
is significantly ahead of the bandwidth-tuned quantum kernels. The gain therefore belongs to the
reference frame and to the kernel formulation, both of which are available
classically and neither of which requires a quantum computer. Once the
invariance groups are matched, the quantum geometries are not better; the
parameter-free ones are not worse either, which is itself informative, since
it locates the ceiling precisely where the algebra says it must be, and the
bandwidth-tuned ones pay a small, partition-dependent price for their extra
hyperparameter.

We draw two lessons. The first is specific: quantum information geometry is a
mathematically exact description of EEG covariances that confers no
decoding benefit over its classical counterparts, and reports to the contrary
should be checked for a frame asymmetry before they are believed. The second is
methodological: the dimension-matched control that this literature has begun to
adopt is necessary but not sufficient. Its natural analogue for geometric
methods, a \emph{metric}-matched control, identical in every respect but the
metric, is what distinguishes a claim about quantum structure from a claim
about preprocessing.

\subsection{Relation to prior work}

Three strands of prior work bear on these results, and our relation to each
is different.

The first is the theory of quantum kernels. Schuld \cite{schuld2021kernel}
showed that supervised quantum models are kernel methods, so their merit
reduces to whether the induced kernel suits the data; Huang \textit{et al}
\cite{huang2021} showed that any advantage is a property of the data as much
as of the kernel; and K\"ubler \textit{et al} \cite{kubler2021} made the
mechanism precise: a quantum kernel helps only where its inductive bias
matches the target. Our results are an instance of that principle in which
the inductive bias is an invariance group. The density-matrix kernels carry
orthogonal invariance, EEG's nuisances generate the congruence group, and the
mismatch alone accounts for their sensor-frame deficit. Bandwidth tuning
\cite{shaydulin2022}, which we apply throughout, and the exponential
concentration of Thanasilp \textit{et al} \cite{thanasilp2024} address a
different failure, one that grows with qubit count; table~\ref{tab:sweep}
shows it is not the failure at work here, since in the reference frame every
density-matrix kernel has more Gram variance at six qubits than at three.

The second is the methodology of QML benchmarking. Bowles \textit{et al}
\cite{bowles2024} found classical models matching or beating quantum ones
across a broad suite and argued for ablations and matched baselines; we adopt
both their entanglement ablation and the dimension-matched control. Our
addition is the \emph{metric}-matched control, which the geometric setting
requires and the feature-map setting does not: a quantum kernel that operates
on an SPD manifold has a classical counterpart on the same manifold, and
comparing it with anything else conflates the metric with the representation.

The third is Riemannian BCI decoding \cite{congedo2017,yger2017}. The
reference-state correction is, algebraically, the re-centring of Zanini
\textit{et al} \cite{zanini2018} and the alignment step of Riemannian
Procrustes analysis \cite{rodrigues2019}, developed there as a
transfer-learning heuristic. What we add is the reason it is not optional: for
a family of kernels with the wrong invariance group it is the difference
between a broken comparison and a fair one, and we prove that it restores
affine invariance exactly rather than approximately. In that light our
findings are compatible with reports of quantum benefit on EEG
\cite{qeegnet2024,eswa2024} only under the interpretation that those benefits
arise from the classical stages of the hybrid pipelines. Our dimension-matched
control speaks to this directly, since a plain linear SVM on the same
PCA-reduced tangent-space features matched or exceeded every circuit kernel,
and our frame result names a second route to the same artefact: a quantum
stage that operates on aligned covariances, compared with a classical baseline
that does not. The pyRiemann-qiskit toolkit \cite{pyriemannqiskit2023}, which
places a quantum classifier after tangent-space projection, is the natural
place for the metric-matched control to be adopted.

\subsection{Limitations}

Five limitations bound our conclusions. (i) We study two datasets but one
binary paradigm, left hand versus right hand; other motor-imagery classes and
other BCI paradigms (P300, SSVEP) are untested. The trial-count concern that a
single dataset would leave open is addressed by section~\ref{sec:crossdataset}.
(ii) Our central claim is an equivalence, and equivalence is bounded rather
than absolute. Two one-sided tests place the quantum kernels and their
classical twins within $\pm\EquivWorstBound$ accuracy of one another across all
\EquivN{} comparisons (table~\ref{tab:equiv}), so we can exclude differences
larger than that, but not smaller ones; under two further partitions of
PhysioNet the bound widens to $\pm\SeedWorstBound$, every departure being in
the twin's favour (table~\ref{tab:seeds}). The TOST implementation is the paired
$t$ form and assumes the per-subject differences are approximately normal; at
$n=\NSubjects$ that is mild, at $n=\BciNSubjects$ it is a real limitation and
the IV-2a bounds should be read as indicative.
(iii) Our invariance argument concerns congruence. Inter-subject and
inter-session shift is dominated by, but not exhausted by, that group;
differences in conditioning, effective rank and spectral shape survive
recentring. Cross-session transfer, where that residual is proportionally
largest, showed no quantum kernel separating from its twin
(section~\ref{sec:crosssession}), but it did so on one dataset of \CsN{}
subjects, which bounds the possibility rather than excluding it.
(iv) The reference state is a Fr\'echet mean estimated from training data, so
the correction is only as good as that estimate; its markedly larger benefit on
IV-2a (\BciFrameMin{} to \BciFrameMax) than on PhysioNet (\PhysFrameMin{} to
\PhysFrameMax) is consistent with this, and implies the correction will be
weaker still in genuinely few-trial calibration settings.
(v) Simulation is noiseless, and although section~\ref{sec:shots} adds shot noise it
models sampling only, not gate error, decoherence or readout bias. All of these
favour the quantum side, so none weakens a negative result; they do mean
hardware performance would be worse, not better. We also survey a specific
family of kernels: deep architectures such as EEGNet \cite{lawhern2018},
and variational or quantum neural-network approaches, are
outside our scope.

\subsection{What would change the conclusion}

Our claim is that the quantum geometries match their classical twins once the
invariance groups agree, so what would overturn it is a setting in which a
quantum divergence separates from a metric-matched classical comparator. We
identified three candidates before running the extended experiments, and two
of them have now been tried.

Non-congruence shift was the first, because it is exactly what our argument
does not cover, and cross-session transfer is where it is proportionally
largest: the congruence component is milder, so residual differences in
conditioning and spectral shape, on which the quantum relative entropy and the
Bures distance weight differently from the affine-invariant metric, carry more
of the shift. Section~\ref{sec:crosssession} reports the outcome: parity, with
no quantum-versus-twin comparison reaching even $p<0.5$. Larger registers
were the third candidate, and the sweep of table~\ref{tab:sweep} shows that
the reference-frame kernels do not concentrate with qubit count. That removes
a concern rather than the parity, since the metric-matched comparison at five
qubits (section~\ref{sec:filterbank}) had already tied.

What remains untried is the regime in which the reference state is poorly
estimated. Our correction depends on a Fr\'echet mean, and the classical
baselines depend on the same estimate; a divergence that degrades more
gracefully under a noisy reference would be a real and practically relevant
advantage, and few-trial calibration is where it would show. The frame
effect's strong dependence on trial count, \PhysFrameMax{} with \NTrials{}
trials per subject against \BciFrameMax{} with \BciNTrials, says that the
reference-state estimate is the binding constraint in that regime, which makes
it the right place to look next.

%% =====================================================================
\section{Conclusion}
%% =====================================================================

Treating EEG covariance matrices as quantum density matrices yields a
principled, hardware-mappable family of quantum kernels that requires no
dimensionality reduction. Evaluated in the sensor frame, these kernels are the
worst pipelines in a controlled benchmark, and appear to be defeated by kernel
concentration. That diagnosis is incomplete: the quantum quantities are
invariant under unitary conjugation while EEG's nuisances act by congruence, so
the comparison was confounded by an invariance mismatch that the classical
baselines do not share. Referring each state to a training-set Fr\'echet mean
restores exact affine invariance, relieves concentration by
\GramGainMin--\GramGainMax$\times$, improves every kernel on both datasets, and
reverses the sign of the classical-versus-quantum comparison.

It would be straightforward to stop there and report a quantum advantage. The
control that prevents this, a classical SPD kernel identical in every respect
but the metric, shows that no quantum kernel exceeds its classical twin in
any of \EquivNSettings{} settings, and that where the two can be told apart at
all, under other partitions of PhysioNet, it is the twin that is ahead. Parity
holds against filter-bank CSP, in
cross-subject transfer, where all nine geometries fall within
\TransferSpread, and in cross-session transfer, the setting we had named in
advance as the most likely to show a real quantum effect. The improvement is
real and useful, and it is entirely classical in origin.

We therefore find no evidence of quantum advantage for motor-imagery decoding,
and we locate the ceiling: matched invariance permits parity and nothing more.
For practitioners, the reference-state correction is worth adopting on its own
terms. For the field, we suggest that a metric-matched control should
accompany the dimension-matched one, since without it these same experiments
support the opposite conclusion. All code and results are released to support
replication.

%% =====================================================================
%% These are iopjournal's structured end-matter commands. All of them except
%% \suppdata are suppressed automatically by the [anonymous] class option, so
%% do NOT hand-anonymise the file for double-anonymous review -- add the
%% option instead.

%% TODO: computing resources, and anyone who helped but is not an author.
\ack{The authors thank the PhysioNet, GigaDB and MOABB maintainers for
curating and openly distributing the datasets used here.}

%% TODO: replace with real funder names and grant numbers, or state that no
%% specific grant was received. Do not leave this as prose if a grant exists;
%% IOP parse this section.
\funding{No specific grant from any funding agency in the public, commercial
or not-for-profit sectors was received for this work.}

%% TODO: CONFIRM. Drafted from the project record, using CRediT terms
%% (https://credit.niso.org). The corresponding author must verify that each
%% listed contribution is accurate and that every author meets authorship
%% criteria; a contribution that does not meet them belongs in \ack instead.
\roles{\textbf{Md. Imtiaj Alam Sajin}: conceptualisation, methodology,
software, formal analysis, investigation, data curation, visualisation,
writing -- original draft. \textbf{Samiul Islam Saif}: software,
investigation. \textbf{Muhammad Hasibur Rashid Chayon}: supervision, writing
-- review and editing.}

\data{The PhysioNet EEG Motor Movement/Imagery Database is publicly available
at \url{https://physionet.org/content/eegmmidb/1.0.0/}, BCI Competition
IV-2a at \url{https://www.bbci.de/competition/iv/}, and the Cho2017 dataset
from GigaDB (\url{https://doi.org/10.5524/100295}); the latter two were
accessed through MOABB. All analysis
code, configuration and result files supporting this study, including the
per-fold scores behind every table and figure, are openly available at
\url{https://github.com/Imtiaj-Sajin/quantaEEG}.}

\section*{Conflicts of interest}
The authors declare no conflicts of interest.

%% No \section*{References} here: under iopjournal the thebibliography
%% environment emits that heading itself, and adding our own printed it twice.
\bibliographystyle{iopart-num}
\bibliography{refs}

\end{document}
